%%%%%%%%%%%%%%%%%%%%%%% file template.tex %%%%%%%%%%%%%%%%%%%%%%%%%
%
% This is a general template file for the LaTeX package SVJour3
% for Springer journals.          Springer Heidelberg 2010/09/16
%
% Copy it to a new file with a new name and use it as the basis
% for your article. Delete % signs as needed.
%
% This template includes a few options for different layouts and
% content for various journals. Please consult a previous issue of
% your journal as needed.
%
%%%%%%%%%%%%%%%%%%%%%%%%%%%%%%%%%%%%%%%%%%%%%%%%%%%%%%%%%%%%%%%%%%%
%
% First comes an example EPS file -- just ignore it and
% proceed on the \documentclass line
% your LaTeX will extract the file if required
% \begin{filecontents*}{example.eps}
%!PS-Adobe-3.0 EPSF-3.0
%%BoundingBox: 19 19 221 221
%%CreationDate: Mon Sep 29 1997
%%Creator: programmed by hand (JK)
%%EndComments
% gsave
% newpath
%   20 20 moveto
%   20 220 lineto
%   220 220 lineto
%   220 20 lineto
% closepath
% 2 setlinewidth
% gsave
%   .4 setgray fill
% grestore
% stroke
% grestore
% \end{filecontents*}
% \RequirePackage{fix-cm}
\documentclass[smallextended,referee,envcountsect]{svjour3}       % onecolumn (second format)
\smartqed  % flush right qed marks, e.g. at end of proof
\usepackage{graphicx}
% \usepackage{textcomp}
% \usepackage{subfigure}
\usepackage[center]{caption}
\usepackage{amsfonts}
% \usepackage{pgfplots}
\usepackage{algorithm}
\usepackage[noend]{algorithmic}
\usepackage{amsmath}

\newtheorem{assumption}{Assumption}

\journalname{To be decided}

\begin{document}

\title{Confluent flow and other approximations to the Pooling Problem}
\titlerunning{Confluent flow and other approximations to the Pooling Problem}

\author{Dag Haugland \and Eligius M.T. Hendrix \and Isak P\`{a}ll Gestsson}

\institute{Dag Haugland, Corresponding author \at 
Department of Informatics, University of Bergen, Norway \email{dag.haugland@ii.uib.no}
\and
Eligius M.T. Hendrix \at
              Department of Computer Architecture, Universidad de M\'alaga and Operations Research and Logistics Group, Wageningen University \email{eligius.hendrix@wur.nl\and Isak P\`{a}ll Gestsson} \at Gassco}\\

\date{Received: date / Accepted: date}

\maketitle

\begin{abstract}
About the computational challenge of the pooling problem and finding approximate solutions by means of confluent and influent flow
\end{abstract}
\keywords{Pooling problem \and complexity \and global optimization}
\subclass{90C26 \and 90C35 \and 90C39 \and 68Q25}


%-------------------------------------------------------------------------------
\section{Introduction}
\label{intro}

Taken from an earlier paper, to be revised:
Traditional network flow models, like the minimum cost network flow problem, rely on considerations of the flow as a homogeneous commodity.
Variations in composition is however relevant in many industrial and logistics applications of the models,
including petroleum refining and pipeline transportation of natural gas.
Crude oils supplied to a refinery have contamination levels depending on their sources of origin.
Major components of natural gas are methane, ethane, butane and propane, but the proportions in which they occur are not equal for all gas wells.
For environmental or technical reasons, requirements to the final composition are often imposed at the sink nodes of the flow network.
Whenever constraints of this kind apply, it is therefore essential for network flow models to recognize not only the total product flow,
but also how the flow composition evolves from network sources to sinks.
In particular, the models must reflect updates of the composition at nodes where differently composed flows are pooled.
This provides a computationally challenging problem referred to as the \emph{pooling problem}.

The pooling problem exhibits strong resemblance with well-studied logistics models like the minimum cost flow problem and the transportation problem.
Another classical related problem is Stigler's diet problem \cite{stigler45,dantzig90},
asking for the minimum cost diet requiring a certain amount of specified nutrients.
The problem is easy to extend to a multi-diet version where groups of people have each their own requirements.
Nutritional contents differ between sources, that is, foods of which diets can be composed,
and the optimal mix must be found for each sink (group of people).
Conceptually, the multi-diet version coincides with the pooling problem in bipartite networks.
Thanks to this simple network structure, with exclusively direct links between sources and sinks, diet problems are nothing but linear programs.
However, when extending the problem to more general network structures, this is no longer the case.
Absence of linear formulations has immediate computational consequences.
While minimum cost flow instances where the network size reaches a six digit number of arcs and nodes can be solved within hours, or even minutes,
exact solution of pooling problem instances with hundreds of nodes appears to be unrealistic.

The pooling problem has been recognized for decades as a challenge where linear programming approaches fail.
Bounds on the flow composition frequently induce bilinear constraints in otherwise linear models,
and require corresponding solution methods.
An approach based on iterated solution of linearized model versions has gained popularity especially in the oil refining industry,
and has been analyzed experimentally by Haverly \cite{haverly78}.
A major conclusion is that methods of this kind may produce suboptimal solutions.
Such observations helped establishing a consensus, already in early days of research in the pooling problem, that global optimality is achievable only in small instances.

Owing to the anticipated intractability of the problem, early research was mainly directed towards heuristic methods \cite{haverly78,baker85,dewitt89},
where iterative linearization is the core idea.
Floudas and Visweswaran \cite{floudas90b} reported the first exact solution algorithm for the pooling problem.
Later, algorithms based on branch-and-bound \cite{audet04,foulds92,misener11,sahinidis05,visweswaran96b},
Lagrangian relaxation \cite{adhya99,almutairi09,bental94}, and integer programming \cite{dey15} have been studied.
The industrial relevance of the problem, notably in petroleum refining \cite{baker85,dewitt89},
the food industry \cite{kallrath05}, and in waste water processing \cite{galan98,misener10b},
has been acknowledged by many authors.
The pooling problem survey by Misener and Floudas \cite{misener09b} has a comprehensive list of references to work in this area.

Standard pooling problems are defined as networks with three node layers, referred to as sources, pools and sinks (sinks), respectively.
Each arc links either a source to a pool or a pool to a sink.
Quality constraints in terms of bounds on the composition are imposed at the sinks.
Each bound relates to the relative content of a flow component, referred to as the component.
Rigorous studies of the computational complexity of the pooling problem have not been conducted until recently.
After NP-hardness was formally proved \cite{alfaki13a}, Dey and Gupte \cite{dey15} showed that the problem is as hard to approximate as the maximum clique problem.

The hardness theorem in \cite{alfaki13a} states that even the class of instances with a unique pool is strongly NP-hard.
The computational complexity is however favorable if there are additional limitations on node cardinalities,
as polynomial time algorithms have been developed for the case of a unique pool and an upper bound on either the number of sinks \cite{haugland15a}
or the number of components \cite{alfaki13a}.
Whether polynomial running-time algorithms for the one-pool instance class exist when the number of sources is bounded, has however been an open question \cite{haugland15a}.

The pooling problem remains NP-hard for networks with only two sources and two sinks,
and only one flow component subject to constraints \cite{haugland15a}.
Whether such instances can be solved in polynomial time when the input data take values within given bounds, has been an open question to date.
In the present work, we answer both of the questions above in the affirmative.
First, we present a new algorithm for the instance class with a unique pool,
and prove that it runs in polynomial time when the source cardinality is fixed. 
Second, we give a dynamic programming algorithm for a rich class of instances with a unique component and only two sources and sinks.
This algorithm is proved to have pseudo-polynomial running time, which implies that its running time is polynomial
for any instance class in which the values of all numerical input data are bounded.
Contributions to the pooling problem literature from the current text are thus answers to two questions left open in existing literature,
and, by virtue of this, an enhanced understanding of the boundaries between tractable and intractable versions of the problem.
 
The remainder of the text is organized as follows:
Mathematical notation and formulation of the pooling problem are provided in Section \ref{sec:problem}.

To be written.

%-----------------------------------------------------------------------
\section{Mathematical Notation and Problem Formulation}\label{sec:problem}

Let $\mathcal{D}=(\mathcal{N},\mathcal{A})$ be a directed graph,
where the node set $\mathcal{N}=S\cup I\cup T$ consists of the disjoint subsets $S$ of \emph{sources}, $I$ of \emph{pools}, and $T$ of \emph{sinks}.
The arc set $\mathcal{A}\subseteq (S\times I)\cup(I\times T)$ links sources with pools and pools with sinks.
Let $K$ be a finite set, and refer to its elements as \emph{flow components} or simply \emph{components}.
In the sequel, elements of vectors and matrices defined over the sets $\mathcal{N}$ and $K$
are identified by a subscript if they correspond to $\mathcal{N}$, and by a superscript if they correspond to $K$.
Complexity analyses in later sections require that network parameters are defined as rational rather than real numbers.
Define the vectors of unit arc costs $\boldsymbol{c}\in\mathbb{Q}^{|\mathcal{A}|}$ and node capacities $\boldsymbol{b}\in\mathbb{Q}_+^{|\mathcal{N}|}$.
For integers $j$ and $n$,
let $\boldsymbol{j}^n$ denote the vector with all its $n$ entries equal to $j$,
and let $\boldsymbol{I}^n$ denote the $(n\times n)$-identity matrix.

At all sources $s\in S$, flow with a given concentration of each component $k\in K$ may enter $\mathcal{D}$, and flow leaving the network at a sink $t\in T$ must satisfy upper concentration bounds.
Define the matrices of \emph{concentration} $\boldsymbol{Q}\in\mathbb{Q}^{|K|\times |S|}$ and \emph{concentration bounds}
$\boldsymbol{U}\in\mathbb{Q}^{|K|\times |T|}$.
For $s\in S$ and $t\in T$, let $\boldsymbol{q}_s=\left(q_s^k\right)_{k\in K}$ and $\boldsymbol{u}_t=\left(u_t^k\right)_{k\in K}$ be the columns of $\boldsymbol{Q}$ and $\boldsymbol{U}$, respectively.
That is, $q_s^k$ is the concentration of component $k$ of the flow entering $\mathcal{D}$ at $s$,
and $u_t^k$ is the upper bound on the concentration of the flow leaving $\mathcal{D}$ at $t$.

For each pool $i\in I$, the set of neighbor sources and the set of neighbor sinks are denoted by
$S_i=\left\{s\in S: (s,i)\in \mathcal{A}\right\}$ and
$T_i=\left\{t\in T: (i,t)\in \mathcal{A}\right\}$, respectively.
For each source $s\in S$,
the set of neighbor pools is denoted by
$I_s=\left\{i\in I: (s,i)\in \mathcal{A}\right\}$, while
$I_t=\left\{i\in I: (i,t)\in \mathcal{A}\right\}$
denotes the set of neighbor pools of sink $t\in T$.
% The arc set $\mathcal{A}$ has a partition $\left(\mathcal{A}_S,\mathcal{A}_T\right)$,
% where $\mathcal{A}_S=\mathcal{A}\cap\left(S\times I\right)$ and $\mathcal{A}_T=\mathcal{A}\cap\left(I\times T\right)$
% denote the sets of source-incident and sink-incident arcs, respectively.

A vector $\boldsymbol{f}\in\mathbb{R}_+^{|\mathcal{A}|}$ is called a \emph{flow} in $\mathcal{D}$ iff it satisfies the capacity and flow conservation constraints:
\begin{eqnarray}
           & \sum_{i\in I_s}f_{si}\leq b_s & s\in S, \label{eq:caps} \\
           & \sum_{s\in S_i}f_{si}\leq b_i & i\in I, \label{eq:capi} \\
           & \sum_{i\in I_t}f_{it}\leq b_t & t\in T, \label{eq:capt} \\
           & \sum_{s\in S_i}f_{si}-\sum_{t\in T_i}f_{it} = 0~~~~ & i\in I. \label{eq:flowcons}
\end{eqnarray}

\noindent
We refer to the set $\mathcal{F}(\mathcal{D},\boldsymbol{b}) = \left\{\boldsymbol{f}\in\mathbb{R}_+^{|\mathcal{A}|}: \mbox{(\ref{eq:caps})--(\ref{eq:flowcons}) are satisfied}\right\}$
as the \emph{flow polytope} of $\mathcal{D}$.
For each $\boldsymbol{f}\in \mathcal{F}(\mathcal{D},\boldsymbol{b})$,
define $\boldsymbol{W}=\left(\boldsymbol{w}_j\right)_{j\in \mathcal{N}}\in\mathbb{R}^{|K|\times |\mathcal{N}|}$ to be a
\emph{concentration matrix} of $\boldsymbol{f}$ iff $\boldsymbol{w}_s=\boldsymbol{q}_s$ ($s\in S$),
\begin{eqnarray*}
  \boldsymbol{w}_i\sum_{t\in T_i}f_{it}=\sum_{s\in S_i}\boldsymbol{q}_sf_{si},\ i\in I,\ \mbox{and} \\
  \boldsymbol{w}_t\sum_{i\in I_t}f_{it}=\sum_{i\in I_t}\boldsymbol{w}_if_{it},\ t\in T.
\end{eqnarray*}
That is, the concentration $\boldsymbol{w}_i=\left(w_i^k\right)_{k\in K}$ of the flow blended at a pool $i$ receiving positive flow,
is a convex combination of the source concentrations $\boldsymbol{q}_s$ ($s\in S$),
with convexity weights $f_{si}/\sum_{t\in T_i}f_{it}$ ($s\in S_i$).
Analogously, for sinks $t\in T$ with positive inflow, the concentration $\boldsymbol{w}_t$ is a convex combination of the pool concentrations $\boldsymbol{w}_i$ ($i\in I_t$).
These conditions on $\boldsymbol{w}$ are referred to as the \emph{linear blending} property, and imply directly
\begin{equation} \label{eq:totqual}
  \sum_{s\in S}\boldsymbol{q}_s\sum_{i\in I_s}f_{si} = \sum_{t\in T}\boldsymbol{w}_t\sum_{i\in I_t}f_{it}.
\end{equation}
At sinks $t\in T$ where $\sum_{i\in I_t}f_{it}>0$, the upper concentration bounds $u_t^k$ apply,
and linear blending yields $\boldsymbol{w}_t=\sum_{i\in I_t}\boldsymbol{w}_if_{it}/\sum_{i\in I_t}f_{it}\leq \boldsymbol{u}_t$.
Thus, $\boldsymbol{f}\in\mathcal{F}\left(\mathcal{D},\boldsymbol{b}\right)$ is said to be a \emph{feasible} flow in
$\mathcal{D}$ iff $\sum_{i\in I_t}\boldsymbol{w}_if_{it}\leq \boldsymbol{u}_t\sum_{i\in I_t}f_{it}$ ($t\in T$)
for any concentration matrix $\boldsymbol{W}$ of $\boldsymbol{f}$.

\begin{problem} \label{prob:pp}
The \textsc{Pooling Problem}: Find a feasible flow $\boldsymbol{f}$ in $\mathcal{D}$ minimizing the total costs $\boldsymbol{c}\cdot\boldsymbol{f}$.
\qed
\end{problem}
Problem \ref{prob:pp}  is traditionally stated in terms of a bilinear formulation in flow and pool concentration variables \cite{haverly78},
often referred to as the P-formulation:

\begin{eqnarray}
\min_{\boldsymbol{f},\boldsymbol{W}} & \left(\sum_{(s,i)\in \mathcal{A}_S}c_{si}f_{si}+\sum_{(i,t)\in \mathcal{A}_T}c_{it}f_{it}\right), & \label{eq:obj} \\
           & \sum_{s\in S_i}q_s^kf_{si}-w_i^k\sum_{t\in T_i}f_{it} = 0~~~~ & i\in I, k\in K, \label{eq:qualcons} \\
           & \sum_{i\in I_t}w_i^kf_{it}-u_t^k\sum_{i\in I_t}f_{it} \leq 0~~~~ & t\in T, k\in K, \label{eq:qualbnd} \\
           & \boldsymbol{f}\in \mathcal{F}(\mathcal{D},\boldsymbol{b}). \label{eq:posflow}
\end{eqnarray}


%-----------------------------------------------------------------
\section{Complexity of the Pooling Problem with a Single Pool} \label{sec:onepool}
In this section, we assume that
$\mathcal{A}=\left(S\times I\right)\cup \left(I\times T\right)$ and $|I|=1$,
which means that all sources and all sinks are connected to a unique pool in the network.
A pool notation thus becomes superfluous, and we let $f_s$ and $f_t$ denote,
respectively, the flow from source $s\in S$ to the pool, and flow from the pool to sink $t\in T$.
An analogous simplified notation is adopted for the arc costs.

\subsection{Problem Formulation with a Set-Valued Variable}
Because a sink $t$ receives flow exclusively from the unique pool, either $f_t=0$,
or the flow entering the pool must make a blend that satisfies the bounds $u_t^k$, $k\in K$.
Hence, we have the following disjunctive constraint for all $t\in T$:
\begin{eqnarray}
 & f_t=0 & ~\mbox{or} \nonumber \\
 & \sum_{s\in S}q_s^kf_s-u_t^k\sum_{s\in S}f_s \leq 0 & ~k\in K. \label{eq:qbnd}
\end{eqnarray}
Considering the set $T^+\subseteq T$ of sinks receiving positive flow as a variable, the problem is formulated as:
\begin{eqnarray}
\min_{T^+,\boldsymbol{f}} & \left(\sum_{s\in S}c_sf_s+\sum_{t\in T}c_tf_t\right), & \label{eq:qobj} \\
           & \boldsymbol{f}\in \mathcal{F}(\mathcal{D},\boldsymbol{b}), & \label{eq:fdom} \\
           & \sum_{s\in S}q_s^kf_s-u_t^k\sum_{s\in S}f_s \leq 0 & k\in K, t\in T^+, \label{eq:t+} \\
           & f_t=0 & t\in T\setminus T^+, \label{eq:t-} \\
           & T^+\subseteq T. & \label{eq:tdom}
\end{eqnarray}
It has been shown \cite{alfaki13a} that the \textsc{Independent Vertex Set} problem is polynomially reducible to problem (\ref{eq:qobj})--(\ref{eq:tdom}),
which hence is (strongly) NP-hard.
By exploiting the same reduction, it was recently shown \cite{dey15} that the problem has no polynomial-time approximation algorithm within $|T|^{1-\epsilon}$, for any $\epsilon>0$,
unless NP-complete problems have randomized polynomial algorithms.

The computational difficulty of problem (\ref{eq:qobj})--(\ref{eq:tdom}) is represented by the set-valued variable $T^+$,
and once $T^+\subseteq T$ is known, the remaining problem (\ref{eq:qobj})--(\ref{eq:t-}) is a linear program (LP).
This observation inspires various solution approaches discussed in the next sections.

For a given $T^+\subseteq T$, let $z(T^+)$ denote the minimum cost in the linear program (\ref{eq:qobj})--(\ref{eq:t-}).

\subsection{Few Terminals} \label{sec:fewt}
If the number $|T|$ of sinks is small, problem (\ref{eq:qobj})--(\ref{eq:tdom}) can be solved by enumerating all $2^{|T|}$ possible $T^+\subseteq T$.
Taking the minimum of the $2^{|T|}$ LP-solutions $z\left(T^+\right)$, yields the optimal solution.
Consequently, the \textsc{Pooling Problem} with a single pool and an upper bound on the number $|T|$ of sinks is solvable in polynomial time \cite{haugland15a}.

\subsection{Few Quality Parameters} \label{sec:fewk}
A small number $|K|$ of components suggests enumeration of all possible pool concentrations rather than enumeration of sink sets.
Consider a component $k\in K$ and a permutation $\pi:\left\{1,\ldots,|T|\right\}\mapsto T$
such that \\ $u_{\pi(1)}^k\leq u_{\pi(2)}^k\leq\cdots\leq u_{\pi(|T|)}^k$.
That is, sink $\pi(1)$ imposes the strictest and sink $\pi(|T|)$ the least strict concentration bound on parameter $k$.
If \\ $u_{\pi(j-1)}^k\sum_{s\in S}f_s < \sum_{s\in S}q_s^kf_s\leq u_{\pi(j)}^k\sum_{s\in S}f_s$,
then the pool concentration satisfies the $k$-bound at sinks $\pi(j),\ldots,\pi(|T|)$,
whereas sinks $\pi(1),\ldots,\pi(j-1)$ must receive zero flow.
Letting $j$ run through all possible values $1,\ldots,|T|+1$, where $j=|T|+1$ indicates that none of the bounds are satisfied,
yields $|T|+1$ possible pairs $\left(\sum_{s\in S}q_s^kf_s-u_{\pi(j)}^k\sum_{s\in S}f_s\leq 0, ~f_{\pi(1)}=\cdots=f_{\pi(j-1)}=0\right)$ of constraints
for this particular $k\in K$.
Combining all $|K|$ components, gives $\left(|T|+1\right)^{|K|}$ compact LPs to solve.
It follows that the \textsc{Pooling Problem} with a single pool and an upper bound on the number $|K|$ of components is also solvable in polynomial time \cite{alfaki13a}.

\subsection{Linearly Ordered Quality Bounds} \label{sec:linord}
Consider instances with total domination in the concentration bounds $\boldsymbol{u}_t$.
This means that the sinks can be ordered such that they appear successively less demanding in all components $k$,
and implies that the permutation $\pi$ discussed in Section \ref{sec:fewk} is independent of $k$.
% That is, for all $j=1,\ldots,|T|-1$, all concentration bounds at sink $\pi(j+1)$ are no stricter than at sink $\pi(j)$.
For the class of such single-pool instances, there exists an algorithm with polynomial running time. This is stated more precisely as follows:
\begin{proposition} \label{prop:linord}
If there exists a bijection $\pi:\left\{1,\ldots,|T|\right\}\mapsto T$ such that $\boldsymbol{u}_{\pi(1)}\leq\cdots\leq\boldsymbol{u}_{\pi(|T|)}$,
then the \textsc{Pooling Problem} with $|I|=1$ is solved in terms of $|T|$ compact linear programs.
\end{proposition}
\begin{proof}
Let $T_j^+=\left\{\pi(j),\ldots,\pi(|T|)\right\}$ ($j=1,\ldots,|T|$) and $T_j^+=\emptyset$ for $j=|T|+1$.
The proposition follows directly by observing that the optimal solution to (\ref{eq:qobj})--(\ref{eq:tdom}) is $\displaystyle{\min_{j=1,\ldots,|T|+1}z\left(T_j^+\right)}$.
\qed
\end{proof}

\begin{remark} \label{rem:pi}
Proposition \ref{prop:linord} also motivates heuristic approaches aimed for
instances where there exists no bijection satisfying the given criterion.
Let $\pi$ be defined such that some measure of the concentration demand at sink $\pi(j)$ is non-increasing with increasing value of $j$.
Then it is reasonable to assume that for some $j=1,\ldots,|T|+1$, the set $T_j^+$ as defined in the above proof will give a small cost $z(T_j^+)$.
The heuristic method hence amounts to solving the $|T|$ associated linear programs.
A natural measure of concentration demand is
$\displaystyle{\sum_{k\in K}\frac{u_t^k}{\max_{s\in S}q_s^k-\min_{s\in S}q_s^k}}$.
It is also possible to apply the indicated heuristic for a set of different measures, e.g.\ by ordering the sinks by increasing value of $u_t^k$,
and thus repeat the heuristic for all $k\in K$.
\qed
\end{remark}

\subsection{Few Sources} \label{sec:fews}
Whether or not there exists an algorithm with polynomial running time for the \textsc{Pooling Problem}
with one pool ($|I|=1$) and an upper bound on the number of sources $|S|$,
has been an open question to date \cite{haugland15a}.
To prepare for an answer in the affirmative, define the map $\mathcal{T}^+:\mathcal{F}(\mathcal{D},\boldsymbol{b})\Rightarrow 2^T$ from the flow polytope to the set of sink subsets,
such that
\[\mathcal{T}^+(\boldsymbol{f})=\left\{\mbox{$t\in T:$ (\ref{eq:qbnd}) is satisfied}\right\}.\]
This means that for a given flow $\boldsymbol{f}$,
$\mathcal{T}^+(\boldsymbol{f})\subseteq T$ consists of \emph{all} sinks $t\in T$ such that $\boldsymbol{f}$ satisfies the concentration constraints of $t$.
A set $T^+\subseteq T$ is thus in the image $\mathcal{T}(\mathcal{D},\boldsymbol{b})$ of $\mathcal{T}^+$,
if and only if there exists an $\boldsymbol{f}\in\mathcal{F}(\mathcal{D},\boldsymbol{b})$ satisfying
all concentration constraints of all $t\in T^+$, while violating at least one concentration constraint of each $t\in T\setminus T^+$.
For instance, under the condition considered in Proposition \ref{prop:linord},
$\mathcal{T}(\mathcal{D},\boldsymbol{b}) = \left\{T_j^+: j=j_0,\ldots,j_1\right\}$ for some $1\leq j_0\leq j_1\leq |T|+1$.

\begin{proposition} \label{prop:z}
$\min\left\{z\left(T^+\right):T^+\subseteq T\right\} = \min\left\{z\left(T^+\right):T^+\in\mathcal{T}\left(\mathcal{D},\boldsymbol{b}\right)\right\}$.
\end{proposition}
\begin{proof}
Because $\boldsymbol{f}=\boldsymbol{0}^{\left|\mathcal{A}\right|}$ satisfies $T$, $T\in\mathcal{T}\left(\mathcal{D},\boldsymbol{b}\right)$,
and the latter minimum exists.
Let $T^+\in 2^T\setminus\mathcal{T}\left(\mathcal{D},\boldsymbol{b}\right)$.
% If $T^+$ is not contained in any set in
% $\mathcal{T}(\mathcal{D},\boldsymbol{b})$,
% then $T^+$ is satisfied only by zero flow, %$\boldsymbol{f}=\boldsymbol{0}^{\left|\mathcal{A}\right|}$,
% implying $z\left(T^+\right)=0$.
Then, any optimal flow in (\ref{eq:qobj})--(\ref{eq:t-}) satisfies some $T^{\ast}\in\mathcal{T}\left(\mathcal{D},\boldsymbol{b}\right)$, where $T^{\ast}\supset T^+$.
It follows that $z\left(T^{\ast}\right)\leq z\left(T^+\right)$.
\qed
\end{proof}

The enumeration algorithm indicated in Section \ref{sec:fewt} computes $z(T^+)$ for \emph{every} $T^+\subseteq T$.
According to Proposition \ref{prop:z}, however,
it suffices to compute $z(T^+)$ for all $T^+\in\mathcal{T}(\mathcal{D},\boldsymbol{b})$.
This is exploited in the algorithm detailed and analyzed in the next section.
Particularly, when $|S|$ is fixed, it is proved that $\left|\mathcal{T}(\mathcal{D},\boldsymbol{b})\right|$ is polynomial in the instance size.
In its turn, this leads to a proof of polynomial running time of the algorithm.

\section{The \textsc{Pooling Problem} with One Pool and Fixed Number of Sources} \label{sec:i1sbnd}

\subsection{Preliminaries} \label{sec:prelim}
For a non-zero flow $\boldsymbol{f}$, let $x_s= \frac{f_s}{\sum_{t\in T}f_t}$ denote the proportion \cite{bental94} of the total flow originating from source $s$.
Without loss of generality \cite[Proposition 2]{haugland15a}, 
we assume that, for some $s_0\in S$, $\boldsymbol{q}_{s_0}=\boldsymbol{0}^{|K|}$.
Letting $S_0=S\setminus\{s_0\}$, we thus have that the proportion vector $\boldsymbol{x}$ takes values on the unit simplex:
\begin{eqnarray}
         & \sum_{s\in S_0}x_s\leq 1, \label{eq:sum1} \\
         & x_s\geq 0 & s\in S_0. \label{eq:pos}
\end{eqnarray}

\noindent
The slack in (\ref{eq:sum1}) equals the flow proportion corresponding to source $s_0$.

Dividing (\ref{eq:qbnd}) by the total flow $\displaystyle{\sum_{s\in S}f_s=\sum_{t\in T}f_t}$ shows that if $T^+\in \mathcal{T}(\mathcal{D},\boldsymbol{b})$,
then there exists an $\boldsymbol{x}$ on the unit simplex such that for all $t\in T^+$:
\begin{equation} \label{eq:t}
\displaystyle{\sum_{s\in S_0}q_s^kx_s\leq u_t^k, k\in K.}
\end{equation}

The hyperplanes satisfying $\sum_{s\in S_0}q_s^kx_s = u_t^k$ for any $t\in T$ and $k\in K$, partition the unit simplex into a set of polytopes.
Across the relative interior of any such polytope, the set $T^+$ of sinks $t$ for which (\ref{eq:t}) is satisfied is invariant.
If all polytopes are available, problem (\ref{eq:qobj})--(\ref{eq:tdom}) is solved by the following operations on each of them:
First, the corresponding set $T^+$ is identified.
This is accomplished by choosing an arbitrary relative interior point $\boldsymbol{x}$, and identifying $T^+$ as $\left\{\mbox{$t\in T$: (\ref{eq:t}) is satisfied}\right\}$.
With $T^+$ fixed, the resulting LP (\ref{eq:qobj})--(\ref{eq:t-}) is solved next.
The best solution hence obtained is also the solution to (\ref{eq:qobj})--(\ref{eq:tdom}).

The outlined approach calls for an algorithm that takes a set of linear equations as input,
and outputs the set of polytopes in the corresponding partition of the unit simplex.
For the case of the \textsc{Pooling Problem} with $|I|=1$,
such an algorithm is illustrated in the next section.
The problem of identifying the partition of the simplex induced by intersecting hyperplanes is analyzed in general in Section \ref{sec:simp}.
The analysis leads to an algorithm presented in Section \ref{sec:alg},
and we prove that its running time is polynomial in the number of hyperplanes when the variable dimension is fixed.
By direct application of this result, we prove in Section \ref{sec:p1pol}
that the \textsc{Pooling Problem} with $|I|=1$ and $|S|$ fixed is also solvable in polynomial time.

\subsection{Illustration} \label{sec:ill}

Consider an instance with sources $S=\left\{s_0,s_1,s_2\right\}$, sinks $T=\left\{t_1,t_2\right\}$, and $|K|=2$.
Let the source concentration matrix be $\boldsymbol{Q} = \left(\boldsymbol{q}_{s_0},\boldsymbol{q}_{s_1},\boldsymbol{q}_{s_2}\right) = \left(
\begin{array}{rrr}
0 &  1 & -2\\
0 & -2 & 1 \\
\end{array} \right)$ and let the concentration bounds be $\boldsymbol{U} = \left(\boldsymbol{u}_{t_1},\boldsymbol{u}_{t_2}\right) = \left(
\begin{array}{rr}
0 & \frac{1}{2} \\
-\frac{1}{2} & 0 \\
\end{array} \right)$.
Thus, $\left|S_0\right|=2$, and (\ref{eq:sum1})--(\ref{eq:pos}) define the two-dimensional unit simplex depicted in Fig.\ \ref{fig:part}.

\begin{figure}[ht]
% \vspace{-1em}
\begin{center}
\includegraphics[scale=0.35]{simpvoorbeeld.pdf}
\vspace{-1em}
\end{center}
\caption{Partition of the unit simplex by hyperplanes}
\label{fig:part}
\end{figure}

Letting $k$ index the first row of $\boldsymbol{Q}$ and $\boldsymbol{U}$, inequality (\ref{eq:t}) becomes (for the case of brevity, write $x_{s_j}=x_j$ ($j=1,2$))
$x_1-2x_2\leq 0$ for $t=t_1$,
and $x_1-2x_2\leq \frac{1}{2}$ for $t=t_2$.
For $t=t_1$,
the inequality is illustrated by a straight line intersecting the boundary of the simplex in Fig.\ \ref{fig:part} at points $(0,0)$ and
$\left(\frac{2}{3},\frac{1}{3}\right)$.
Analogously, for $t=t_2$, the intersection points are 
$\left(\frac{1}{2},0\right)$ and
$\left(\frac{5}{6},\frac{1}{6}\right)$.
Inequalities
$-2x_1+x_2\leq \frac{-1}{2}$
and $-2x_1+x_2\leq 0$,
corresponding to the second row of $\boldsymbol{Q}$ and $\boldsymbol{U}$,
give the remaining two lines intersecting the simplex.

The hyperplanes consisting of all $\boldsymbol{x}$-vectors satisfying one of the constraints (\ref{eq:t}) with equality,
thus partition the unit simplex into polytopes denoted a, b, c, d, e and f.
We observe that polytope d corresponds to $T^+=\{t_1,t_2\}$, polytopes b, c and e correspond to $T^+=\{t_2\}$, and polytopes a and f correspond to $T^+=\emptyset$.
Thus, $\mathcal{T}(\mathcal{D},\boldsymbol{b})=\left\{T,\{t_2\},\emptyset\right\}$, which means that the LP (\ref{eq:qobj})--(\ref{eq:t-}) does not have to be solved for $T^+=\{t_1\}$.

In Fig.\ \ref{fig:part}, there are 10 points appearing as a vertex in at least one of the polytopes.
These are (1/3, 1/6) and 9 points on the boundary of the simplex.
Each vertex point is identified by the solution to a system of linear equations.
A pair of vertex points is further an edge of some polytope if and only if
there exist $n-1$ distinct hyperplanes corresponding to (\ref{eq:sum1})--(\ref{eq:t}) intersecting both points,
and no such hyperplane separating the points.
We observe that in Fig.\ \ref{fig:part}, $[(0,0),(1/3,1/6)]$ is an edge of c and e.
Due to the line separating $(5/6,1/6)$ and $(1/2,1/2)$, $[(5/6,1/6),(1/2,1/2)]$ is not an edge of any polytope.
Neither is $[(1/3,1/6),(1/3,2/3)]$, since there is no line corresponding to (\ref{eq:sum1})--(\ref{eq:t}) intersecting $(1/3,1/6)$ and $(1/3,2/3)$.

In addition to the polytopes in the partition, also their proper faces are of interest.
In cases where two or more of the hyperplanes corresponding to (\ref{eq:t}) coincide,
some $T^+\in\mathcal{T}(\mathcal{D},\boldsymbol{b})$ can map to a polytope of dimension smaller than $|S_0|$.
The algorithm to be developed next therefore produces a list of all faces of polytopes in the partition of the simplex.


\subsection{Polyhedral Partitions of the Unit Simplex} \label{sec:simp}

\subsubsection{Definitions} \label{sec:def}

Let $P\subset\mathbb{R}^n$ denote a polytope and $V(P)$ its vertex set.
For a finite set $V\subset\mathbb{R}^n$, let $P(V)$ denote the polytope $\mbox{conv}V$.
A polytope of dimension $n$ is referred to as an $n$-polytope, and analogous nomenclature is used for other geometrical structures in vector spaces.

Define the index sets $M=\{1,\ldots,m+n+1\}$ and $N=\{1,\ldots,n\}$, and consider a matrix $\boldsymbol{A}\in\mathbb{Q}^{m\times n}$ and a vector $\boldsymbol{d}\in\mathbb{Q}^m$.
Associate with $(\boldsymbol{A},\boldsymbol{d})$ the extension $(\bar{\boldsymbol{A}},\bar{\boldsymbol{d}})$, where
\[
\bar{\boldsymbol{A}} = \left(
\begin{array}{c}
\boldsymbol{A} \\
\boldsymbol{I}^n \\
% \left(\boldsymbol{1}^n\right)^T
1\ \cdots\ 1
\end{array} \right) = \left(a_{ij}\right)_{i\in M, j\in N},\ 
\bar{\boldsymbol{d}} = \left(
\begin{array}{c}
\boldsymbol{d} \\
\boldsymbol{0}^n \\
1
\end{array} \right) = \left(d_i\right)_{i\in M}.
\]

\noindent
Observe that the last $n+1$ rows of $(\bar{\boldsymbol{A}},\bar{\boldsymbol{d}})$ contain the coefficients of the linear inequalities 
$\sum_{j=1}^nx_j \leq 1$ and $x_j \geq 0$ ($j=1,\ldots,n$)
defining the unit $n$-simplex $\Delta^n$.

\begin{definition}\label{def:vc}
A pair $\left(M^<,M^>\right)$ of sets $M^<\subseteq M$ and $M^>\subseteq M$ of row indices of $(\bar{\boldsymbol{A}},\bar{\boldsymbol{d}})$,
is referred to as a \emph{valid cover} iff
\begin{enumerate}
  \item $M^<\cup M^>=M$ (the rows of $(\bar{\boldsymbol{A}},\bar{\boldsymbol{d}})$ are covered), \label{eq:vc1}
  \item $m+1,\ldots,m+n\in M^>$ ($M^>$ covers the rows corresponding to \\ $x_1,\ldots,x_n\geq 0$), and \label{eq:vc2}
  \item $m+n+1\in M^<$ ($M^<$ covers the row corresponding to $\sum_{j=1}^nx_j\leq 1$). \label{eq:vc3}
\qed
\end{enumerate}
\end{definition}

\noindent
It is not required that $M^<\cap M^>$ is empty.
Conditions \ref{eq:vc2}--\ref{eq:vc3} of Definition \ref{def:vc} imply that
the polytope consisting of all $\boldsymbol{x}\in\mathbb{R}^n$ satisfying
\begin{eqnarray}
 & \sum_{j=1}^na_{ij}x_j \leq d_i & i\in M^<, \label{eq:leq} \\
 & \sum_{j=1}^na_{ij}x_j \geq d_i & i\in M^>, \label{eq:geq}
\end{eqnarray}
\noindent
is contained in $\Delta^n$.
Let $\mathcal{P}^r(\boldsymbol{A},\boldsymbol{d})$ be the set of $r$-polytopes ($r=-1,0,1,\ldots,n$) that correspond to valid covers,
and let $\mathcal{V}^r(\boldsymbol{A},\boldsymbol{d})$ be the set of vertex sets of these polytopes.
Define 
\[ \mathcal{P}(\boldsymbol{A},\boldsymbol{d}) = \mathcal{P}^{-1}(\boldsymbol{A},\boldsymbol{d})\cup\cdots\cup\mathcal{P}^n(\boldsymbol{A},\boldsymbol{d}) ~~\mbox{and} \]
\[\mathcal{V}(\boldsymbol{A},\boldsymbol{d}) = \mathcal{V}^{-1}(\boldsymbol{A},\boldsymbol{d})\cup\cdots\cup\mathcal{V}^n(\boldsymbol{A},\boldsymbol{d}). \]
Here, the convention $\dim\emptyset=-1$ is used,
i.e., $\mathcal{P}^{-1}(\boldsymbol{A},\boldsymbol{d})=\emptyset$.
All valid covers $(M^<,M^>)$ where
$m+1,\ldots,m+n\in M^<$ and $m+n+1\in M^>$ correspond to $\emptyset$, but also others may do so.

Corresponding to a polytope $P\in\mathcal{P}(\boldsymbol{A},\boldsymbol{d})$, there is at least one valid cover.
In degenerate cases, like $P$ containing only the intersection point of three lines in the plane, there are more than one valid cover corresponding to $P$.
A valid cover $\left(M^<,M^>\right)$ is said to be \emph{maximal} iff, for some $P\in\mathcal{P}(\boldsymbol{A},\boldsymbol{d})$,
\[ M^<=\left\{i\in M: \sum_{j=1} ^na_{ij}x_j\leq d_i ~\forall x\in P\right\} \mbox{and} \]
\[ M^>=\left\{i\in M: \sum_{j=1} ^na_{ij}x_j\geq d_i ~\forall x\in P\right\}.\]
That is, $\left(M^<,M^>\right)$ is maximal if and only if $M^<$ and $M^>$ contain the indices of \emph{all} $i\in M$
such that $\sum_{j=1}^na_{ij}x_j\leq d_i$ and $\sum_{j=1}^na_{ij}x_j\geq d_i$, respectively,
for all $\boldsymbol{x}\in P$.
Consequently, there is a one-to-one correspondence between polytopes in $\mathcal{P}(\boldsymbol{A},\boldsymbol{d})$ and maximal valid covers.
In the example of a point intersected by three lines, the indices of all three lines occur in both $M^<$ and $M^>$ in the corresponding maximal valid cover.

\begin{definition}[{\cite[Ch.\ IV, \S 1.2]{aleksandrov56}},{\cite[Definition 2.1.5]{deloera10}}] \label{def:polcomplex}
A finite set $\mathcal{C}$ of polyhedra in $\mathbb{R}^n$ is a \emph{polyhedral complex} iff
\begin{enumerate}
  \item $F\in\mathcal{C}$ if $F$ is a face of a polyhedron $P\in\mathcal{C}$, and
  \item $P_0\cap P_1$ is a face of both $P_0$ and $P_1$ if $P_0,P_1\in\mathcal{C}$.
\qed
\end{enumerate}
\end{definition}

To prove that the size of
$\mathcal{V}(\boldsymbol{A},\boldsymbol{d})$ and
$\mathcal{P}(\boldsymbol{A},\boldsymbol{d})$
is polynomial in $m$,
the concept of a polyhedral complex is instrumental.
% This result is crucial for proving that the \textsc{Pooling Problem} with $|I|=1$ and $|S|$ fixed has a polynomial time algorithm.

\subsubsection{Preliminary Results}

\begin{proposition} \label{prop:polcomplex}
$\mathcal{P}(\boldsymbol{A},\boldsymbol{d})$ is a polyhedral complex.
\end{proposition}

\begin{proof}
To prove that $\mathcal{P}(\boldsymbol{A},\boldsymbol{d})$ satisfies the first condition of Definition \ref{def:polcomplex},
consider any $P\in \mathcal{P}^r(\boldsymbol{A},\boldsymbol{d})$,
and let $F$ be an $r'$-face of $P$, where $0\leq r'<r\leq n$.
Then there exists a valid cover
$\left(M^<,M^>\right)$ for which $P$ is the solution set of (\ref{eq:leq})--(\ref{eq:geq}).
Because $F$ is an $r'$-face of $P$, there also exist $r-r'$ linearly independent inequalities in
(\ref{eq:leq})--(\ref{eq:geq}),
that are satisfied with equality for all $\boldsymbol{x}\in F$,
and satisfied with strict inequality for all $\boldsymbol{x}\in P\setminus F\neq\emptyset$.
Let $L^<\subseteq M^<$ and $L^>\subseteq M^>$ ($\left|L^<\right|+\left|L^>\right|=r-r'$), be the indices of the inequalities meeting these criteria.
Since $F$ is the polytope corresponding to the valid cover $\left(M^<\cup L^>,M^>\cup L^<\right)$,
it follows that $F\in \mathcal{P}^{r'}(\boldsymbol{A},\boldsymbol{d})\subset\mathcal{P}(\boldsymbol{A},\boldsymbol{d})$.

To prove the second condition, consider any $P_0,P_1\in\mathcal{P}(\boldsymbol{A},\boldsymbol{d})$
with corresponding maximal covers
$\left(M_0^<,M_0^>\right)$ and
$\left(M_1^<,M_1^>\right)$, respectively.
Thus, $P_0\cap P_1$ consists of all $\boldsymbol{x}\in \mathbb{R}^n$ satisfying
$\sum_{j=1}^na_{ij}x_j\leq d_i$ $\left(i\in M_0^<\cup M_1^<\right)$ and $\sum_{j=1}^na_{ij}x_j\geq d_i$ $\left(i\in M_0^>\cup M_1^>\right)$.
Let $M^==\left(M_0^<\cap M_1^>\right)\cup\left(M_1^<\cap M_0^>\right)$ and
$H=\left\{\boldsymbol{x}\in \mathbb{R}^n:\sum_{j=1}^na_{ij}x_j=d_i, i\in M^=\right\}$.
For $k=0,1$, it follows that 
\[
  P_k\cap H=\left\{\boldsymbol{x}\in\mathbb{R}^n: \sum_{j=1}^na_{ij}x_j\leq d_i \left(i\in \widehat{M}_k^<\right),\sum_{j=1}^na_{ij}x_j\geq d_i \left(i\in \widehat{M}_k^>\right)\right\},
\]
\noindent
where
\[ \widehat{M}_k^<=M_k^<\cup M^==M_k^<\cup\left(M_k^<\cap M_{1-k}^>\right)\cup\left(M_{1-k}^<\cap M_k^>\right)= M_k^<\cup M_{1-k}^<.\]
% $= M_k^<\cup\left(M_{1-k}^<\cap M_k^>\right)
% = M_k^<\cup\left(M_k^<\cap M_{1-k}^<\right)\cup\left(M_{1-k}^<\cap M_k^>\right)

The last equality is obtained by exploiting $M_k^<\cup M_k^>=M$,
which is true since $\left(M_k^<,M_k^>\right)$ is a valid cover.
Analogously, $\widehat{M}_k^>=M_k^>\cup M_{1-k}^>$,
and hence, $P_k\cap H = P_0\cap P_1$.
% for $k=1,2$,
% $P_k=\left\{x\in \mathbb{R}^n: \forall\left(i\in M_k^<\right)\sum_{j=1}^na_{ij}x_j\leq d_i, \forall\left(i\in M_k^>\right)\sum_{j=1}^na_{ij}x_j\geq d_i\right\}$, and
% $P_0\cap P_1 = H\cap
% \left\{x\in \mathbb{R}^n: \forall\left(i\in M_0^<\cup M_1^<\right)\sum_{j=1}^na_{ij}x_j\leq d_i, \forall\left(i\in M_0^>\cup M_1^>\right)\sum_{j=1}^na_{ij}x_j\geq d_i\right\}$.
Because % all $\boldsymbol{x}\in P_k$ ($k=1,2$) satisfy
% $\sum_{j=1}^na_{ij}x_j\leq b_i$ for all $i\in M^=\cap M_k^<$ and
% $\sum_{j=1}^na_{ij}x_j\geq b_i$ for all $i\in M^=\cap M_k^>$,
$P_k$ is entirely contained in one of the two closed half-spaces bounded by $H$,
it follows that $P_0\cap P_1$ is a face of $P_k$. % which completes the proof.
\qed
\end{proof}

\begin{proposition} \label{prop:cover}
For all $r=0,1,\ldots,n$, $\displaystyle{\left|\mathcal{V}^r(\boldsymbol{A},\boldsymbol{d})\right|\leq {\left|\mathcal{V}^0(\boldsymbol{A},\boldsymbol{d})\right|\choose r+1}}$.
\end{proposition}
\begin{proof}
Let $p_r=\left|\mathcal{V}^r(\boldsymbol{A},\boldsymbol{d})\right|$, and $\mathcal{V}^r(\boldsymbol{A},\boldsymbol{d})=\left\{V_1,\ldots,V_{p_r}\right\}$.
For all $k=1,\ldots,p_r$, $P\left(V_k\right)$ has a simplicial decomposition \cite{edmonds70}, i.e. there exist $t_k\geq 1$ subsets
$U_{k1},\ldots,U_{kt_k}\subseteq V_k$ with $\left|U_{k1}\right|=\cdots=\left|U_{kt_k}\right|=r+1$ such that  
$P\left(V_k\right)$ is 
covered by the set of $r$-simplices $P\left(U_{k1}\right),\ldots,P\left(U_{kt_k}\right)$.
By Proposition \ref{prop:polcomplex}, $U_{k1}\ldots,U_{kt_k}$ consist of vertices from $\mathcal{V}^0(\boldsymbol{A},\boldsymbol{d})$.

Consider two polytopes $P\left(V_k\right),P\left(V_{\ell}\right)\in \mathcal{P}^r(\boldsymbol{A},\boldsymbol{d})$, where $V_k\neq V_{\ell}$.
% and their simplicial decompositions $U_{k1},\ldots,U_{kt_k}$ and $U_{\ell1},\ldots,U_{\ell t_{\ell}}$, respectively.
Assume there exist some $j[k]\in \{1,\ldots,t_k\}$ and $j[\ell]\in \{1,\ldots,t_{\ell}\}$ such that \\ $U_{k,j[k]}=U_{\ell,j[\ell]}=:U$.
Because Proposition \ref{prop:polcomplex} implies that $P\left(V_k\right)\cap P\left(V_{\ell}\right)$ is a face of both
$P\left(V_k\right)$ and $P\left(V_{\ell}\right)$, we have
\[r=\dim{P(U)}\leq\dim{\left(P\left(V_k\right)\cap P\left(V_{\ell}\right)\right)}\leq\dim{P\left(V_k\right)}=\dim{P\left(V_{\ell}\right)}=r.\]
Hence, $P\left(V_k\right)\cap P\left(V_{\ell}\right)$ equals $P\left(V_k\right)$ and $P\left(V_{\ell}\right)$, contradicting $V_k\neq V_{\ell}$.
It follows that $U_{11},\ldots,U_{p_r1}$ are distinct sets consisting of $r+1$ points in $\mathcal{V}^0(\boldsymbol{A},\boldsymbol{d})$.
% It follows that for the vertex sets $V_1,\ldots,V_{p_r}$ in $\mathcal{V}^r(\boldsymbol{A},\boldsymbol{d})$,
% there exist distinct subsets $U_{1,j[1]},\ldots,U_{p_r,j[p_r]}$, each consisting of $r+1$ elements from $\mathcal{V}^0(\boldsymbol{A},\boldsymbol{d})$.
As there exist exactly $\displaystyle{\left|\mathcal{V}^0(\boldsymbol{A},\boldsymbol{d})\right|\choose r+1}$ such sets, the proof is complete.
\qed
\end{proof}

\subsection{Enumerating all Elements in the Polyhedral Complex} \label{sec:alg}

Enumeration of all faces of a polyhedron is an interesting problem that has received considerable attention \cite{murty95,fukuda97} in computational geometry studies.
In this section, an algorithm for the following more general problem is analyzed:
\begin{problem} \label{prob:part}
\textsc{Enumeration of a Polyhedral Complex}: \\ Given $\boldsymbol{A}\in\mathbb{Q}^{m\times n}$ and $\boldsymbol{d}\in\mathbb{Q}^m$,
generate a list of all vertex sets of polytopes in the polyhedral complex $\mathcal{P}(\boldsymbol{A},\boldsymbol{d})$.
\qed
\end{problem}

Below, it is formally proved that Algorithm \ref{alg:fpol} produces
a list of all polytopes in the polyhedral complex $\mathcal{P}(\boldsymbol{A},\boldsymbol{d})$.
Each such polytope, $P\in \mathcal{P}(\boldsymbol{A},\boldsymbol{d})$, is represented by its vertex set $V(P)$.
First, the list $\mathcal{V}^0(\boldsymbol{A},\boldsymbol{d})$ consisting of singletons is generated by selecting sets
$L\subseteq M$ indexing $n$ linearly independent rows of $\bar{\boldsymbol{A}}$.
Denote the indexed rows of $\left(\bar{\boldsymbol{A}},\bar{\boldsymbol{d}}\right)$ by $\left(\bar{\boldsymbol{A}}_L,\bar{\boldsymbol{d}}_L\right)$.
Solutions to the corresponding linear systems that are contained on the unit simplex are stored in $\mathcal{V}^0(\boldsymbol{A},\boldsymbol{d})$, whereas others are discarded.

Next, for increasing values of $r$, all vertex sets in $\mathcal{V}^{r+1}(\boldsymbol{A},\boldsymbol{d})$ are found by extending vertex sets of $\mathcal{V}^r(\boldsymbol{A},\boldsymbol{d})$.
% This is done by considering all sets $V$ constructed as follows.
For each $U\in\mathcal{V}^r(\boldsymbol{A},\boldsymbol{d})$ and $W\in\mathcal{V}^0(\boldsymbol{A},\boldsymbol{d})$,
let $V$ consist of the vertices in $\mathcal{V}^0(\boldsymbol{A},\boldsymbol{d})$ that satisfy all inequalities in (\ref{eq:leq})--(\ref{eq:geq}) satisfied by 
$U\cup W$.
Thus, $V\supseteq U\cup W$.
% If any of the hyperplanes $\sum_{j=1}^na_{ij}x_j=d_i$ ($i\in M$) separates $V^r$ from $V^0$, then $V$ contains vertices of different $(r+1)$-polytopes in the polyhedral complex,
% and is not considered further.
If the inequalities satisfied by $V$ constitute a valid cover,
and $\dim{P(V)}=r+1$,
$V$ is added to $\mathcal{V}^r(\boldsymbol{A},\boldsymbol{d})$.
% we search the current list $\mathcal{V}^{r+1}(\boldsymbol{A},\boldsymbol{d})$ for a set $W$ of
% vertices that neither can be separated from $V$.
% If one such set is found, it is extended by the vertices in $V$, reflecting the fact that the extended set consists of vertices of an $(r+1)$-polytope in the polyhedral complex.
% Otherwise, $V$ is added to the list as a new item.

\begin{algorithm} \caption{\texttt{findPolyhedralComplex}($\boldsymbol{A}$, $\boldsymbol{d}$)} \label{alg:fpol}
\begin{algorithmic} 
  \STATE $\mathcal{V}^0\leftarrow\emptyset$
  \FOR{all sets $L\subset M$ indexing $n$ linearly independent rows of $\bar{\boldsymbol{A}}$}
    \STATE $\boldsymbol{x}\leftarrow \bar{\boldsymbol{A}}_{L}^{-1}\bar{\boldsymbol{d}}_{L}$, $V\leftarrow\{\boldsymbol{x}\}$
    \IF{$\boldsymbol{x}\in \Delta^n$ \AND $V\not\in\mathcal{V}^0$}
      \STATE $\mathcal{V}^0\leftarrow\mathcal{V}^0\cup\{V\}$,
      \STATE $M_V^<\leftarrow \left\{i\in M: \sum_{j=1}^na_{ij}x_j\leq d_i\right\}$,
             $M_V^>\leftarrow \left\{i\in M: \sum_{j=1}^na_{ij}x_j\geq d_i\right\}$
    \ENDIF
  \ENDFOR
  \COMMENT{Now $\mathcal{V}^0$ contains all singular point sets}
  \FOR{$r\leftarrow 0,\ldots,n-1$}
    \STATE $\mathcal{V}^{r+1}\leftarrow\emptyset$  \COMMENT{No $(r+1)$-polytope found yet}
    \FOR{all $U\in\mathcal{V}^r$ and $W\in\mathcal{V}^0$ where $W\not\subseteq U$}
      \STATE $M^<\leftarrow M_{U}^<\cap M_{W}^<$, $M^>\leftarrow M_{U}^>\cap M_{W}^>$
      \IF[Valid cover, no $\left(\bar{\boldsymbol{A}},\bar{\boldsymbol{d}}\right)$-hyperplane separates $W$ from $U$]{$M^<\cup M^>=M$}
        \STATE $V\leftarrow\left\{\boldsymbol{x}\in\mathcal{V}^0:
               \sum_{j=1}^na_{ij}x_j\leq b_i (i\in M^<), \sum_{j=1}^na_{ij}x_j\geq b_i (i\in M^>)\right\}$
        \IF[consider $V$ only if $V\in\mathcal{V}^{r+1}\left(\boldsymbol{A},\boldsymbol{d}\right)$]{$\dim{P(V)}=r+1$}
            \STATE $\mathcal{V}^{r+1} \leftarrow \mathcal{V}^{r+1}\cup\{V\}$, $M_V^<\leftarrow M^<$, $M_V^>\leftarrow M^>$ \COMMENT{an $(r+1)$-polytope is found}
%           \STATE \COMMENT{include the vertex set of an $(r+1)$-polytope in the polyhedral complex}
        \ENDIF
      \ENDIF
    \ENDFOR
  \ENDFOR
  \RETURN{$\left(\mathcal{V}^0,\ldots,\mathcal{V}^n\right)$}
\end{algorithmic} 
\end{algorithm} 

For an illustration of Algorithm \ref{alg:fpol}, assume that $\boldsymbol{A}$ and $\boldsymbol{d}$ correspond to the example depicted in Fig.\ \ref{fig:part}.
Consider the iterations of the second and third loops of the algorithm, in which $r$ is assigned the value 0,
while $U$ and $W$ are assigned values $\{(1/3,1/6)\}$ and $\{(1/2,1/2)\}$, respectively.
There is no hyperplane separating the two points, and a valid cover is identified.
Only points on the line segment $[(1/3,1/6),(1/2,1/2)]$ satisfy all inequalities satisfied by $U\cup W$, and thus 
$V$ becomes $\{(1/3,1/6), (1/2,1/2)\}$.
Since also $\dim{P(V)}=1$, $V$ is added to $\mathcal{V}^1$.
When $U=\{(5/6,1/6)\}$ and $W=\{(1/2,1/2)\}$, however, the hyperplane intersecting (0,0) and (2/3,1/3) separates $U$ and $W$.
As no valid cover is found, the condition of the first \texttt{if}-statement inside the last loop does not hold.
When $U=\{(1/3,1/6)\}$ and $W=\{(1/3,2/3)\}$, the condition holds, and the value $\{(0,0),(1/3,1/6),(1/2,1/2),(1/3,2/3)\}$ is assigned to $V$.
However, $\dim{P(V)}=2$, and $V$ is not added to $\mathcal{V}^1$.
In the next iteration of the loop on $r$, $U$ will in some iteration of the innermost loop be assigned the value $\{(1/3,1/6), (1/2,1/2)\}$ while $W=\{(1/3,2/3)\}$.
Then, $V$ again becomes $\{(0,0),(1/3,1/6),(1/2,1/2),(1/3,2/3)\}$, which finally is added to $\mathcal{V}^2$.

% and the condition of the \texttt{if}-statement with label 1 holds.
% The same is true for the \texttt{if}-statement with label 2, because 
% Assuming that there exists no $W\in\mathcal{V}^2$ satisfying the condition of the last \texttt{if}-statement,
% $V$ is added to $\mathcal{V}^2$.
% In a later iteration of the third loop,
% $V^1$ and $W$ are assigned values $[(0,0),(1/3,1/6)]$ and $(1,3,1/3)$, respectively.
% Now, $V$ becomes $[(0,0), (1/3,1/6), (1/3,2/3)]$, and $\left(M_V^<,M_V^>\right)$ is again a valid cover.
% However, the triangle $[(1/3,1/6), (1/2,1/2), (1/3,2/3)]$ added to $\mathcal{V}^2$ in the iteration discussed first,
% now satisfies the condition of the last \texttt{if}-statement.
% Consequently, $W=[(1/3,1/6), (1/2,1/2), (1/3,2/3)]$, and the quadrangle $V\cup W=[(0,0), (1/3,1/6), (1/2,1/2), (1/3,2/3)]$ replaces $W$ in $\mathcal{V}^2$.

% Consider the iterations of the second and third loops of the algorithm, in which $r$ is assigned the value 1,
% while $V^1$ and $V^0$ are assigned values $[(1/3,1/6),(1/2,1/2)]$ and $(1,3,1/3)$, respectively.
% Subsequently, $V$ becomes the triangle $[(1/3,1/6), (1/2,1/2), (1/3,2/3)]$, $\left(M_V^<,M_V^>\right)$ is a valid cover,
% and the condition of the \texttt{if}-statement with label 1 holds.
% The same is true for the \texttt{if}-statement with label 2, because 
% Assuming that there exists no $W\in\mathcal{V}^2$ satisfying the condition of the last \texttt{if}-statement,
% $V$ is added to $\mathcal{V}^2$.
% In a later iteration of the third loop,
% $V^1$ and $V^0$ are assigned values $[(0,0),(1/3,1/6)]$ and $(1,3,1/3)$, respectively.
% Now, $V$ becomes $[(0,0), (1/3,1/6), (1/3,2/3)]$, and $\left(M_V^<,M_V^>\right)$ is again a valid cover.
% However, the triangle $[(1/3,1/6), (1/2,1/2), (1/3,2/3)]$ added to $\mathcal{V}^2$ in the iteration discussed first,
% now satisfies the condition of the last \texttt{if}-statement.
% Consequently, $W=[(1/3,1/6), (1/2,1/2), (1/3,2/3)]$, and the quadrangle $V\cup W=[(0,0), (1/3,1/6), (1/2,1/2), (1/3,2/3)]$ replaces $W$ in $\mathcal{V}^2$.

The following result shows that Algorithm \ref{alg:fpol} generates the desired lists of vertex sets of polytopes in the polyhedral complex.

\begin{lemma} \label{lem:alg}
  Let $r\in\{0,1,\ldots,n\}$.
  After $r$ iterations of the second \texttt{for}-loop of Algorithm \ref{alg:fpol},
  $\mathcal{V}^\ell=\mathcal{V}^\ell(\boldsymbol{A},\boldsymbol{d})$ for all $\ell\in\{0,1,\ldots,r\}$.
  For all $V\in \mathcal{V}^\ell$, $\left(M_V^<,M_V^>\right)$ is the maximal valid cover corresponding to $P(V)$.
\end{lemma}
\begin{proof}
The proof is by induction in $r$.

Vertex set $\mathcal{V}^0(\boldsymbol{A},\boldsymbol{d})$ consists of all singleton sets $\{\boldsymbol{x}\}$ such that $\boldsymbol{x}\in\Delta^n$ is the unique point satisfying $n$
of the equations in the system $\bar{\boldsymbol{A}}\boldsymbol{x}=\bar{\boldsymbol{d}}$.
This means that the first loop generates $\mathcal{V}^0(\boldsymbol{A},\boldsymbol{d})$, and the claims hold for $r=0$.

To prove the induction step, choose $r\in\{0,1,\ldots,n-1\}$, and assume the claims hold for $r$.
The induction hypothesis says that 
$\left(M_{U}^<,M_{U}^>\right)$ and $\left(M_{W}^<,M_{W}^>\right)$ are the maximal valid covers corresponding to
$P\left(U\right)\in\mathcal{P}^r\left(\boldsymbol{A},\boldsymbol{d}\right)$ and
$P\left(W\right)\in\mathcal{P}^0\left(\boldsymbol{A},\boldsymbol{d}\right)$, respectively.
If $\left(M^<,M^>\right)=\left(M_{U}^<\cap M_{W}^<, M_{U}^>\cap M_{W}^>\right)$ is valid,
it is the maximal valid cover corresponding to a polytope $P\in\mathcal{P}\left(\boldsymbol{A},\boldsymbol{d}\right)$ containing $U\cup W$.
Because the value assigned to $V$ is the set of vertices in $\mathcal{V}^0\left(\boldsymbol{A},\boldsymbol{d}\right)$ satisfying (\ref{eq:leq})--(\ref{eq:geq}),
it follows that $V$ is the vertex set of $P$.
If $\dim{P}=r+1$, $V$ is added to $\mathcal{V}^{r+1}$, and the corresponding maximal valid cover is assigned to $\left(M_V^<,M_V^>\right)$.
Consequently, every iteration of the innermost \texttt{for}-loop leaves $\mathcal{V}^{r+1}$ unchanged, or extends it by the vertex set of an $(r+1)$-polytope.

To complete the proof, we must show that every vertex set
$V\in\mathcal{V}^{r+1}\left(\boldsymbol{A},\boldsymbol{d}\right)$ is added to $\mathcal{V}^{r+1}$.
By the selection of $U$ and $W$, it follows from Proposition \ref{prop:polcomplex}
% that $\left(P\left(V^r\right),W\right)$ enumerates all pairs $(F,\{\boldsymbol{v}\})$ consisting of a facet $F$ and a vertex $\boldsymbol{v}\not\in F$ of $P$.
that in some iteration, $U$ is the vertex set of a facet of $P(V)$, while $W\not\subseteq U$ is a vertex of $P(V)$.
Because $P(V)$ is the unique $(r+1)$-polytope in $\mathcal{P}^{r+1}\left(\boldsymbol{A},\boldsymbol{d}\right)$ containing $U$ and $W$,
it follows that $V$ will be included in $\mathcal{V}^{r+1}$.
\qed
\end{proof}


\begin{proposition} \label{prop:ptop}
  Algorithm \ref{alg:fpol} solves Problem \ref{prob:part}, and its running time is polynomial in $m$.
\end{proposition}
\begin{proof}
Correctness of the output  of Algorithm \ref{alg:fpol} follows directly from Lemma \ref{lem:alg}.
The first \texttt{for}-loop amounts to solving ${m+n+1\choose n}$ systems of $n$ linear equations, and checking $2(m+1)$ linear inequalities.
Consequently, the loop runs in time bounded by a polynomial in $m$.

By Lemma \ref{lem:alg}, the number of sets $V$ considered in iteration $r$ of the second \texttt{for}-loop is no more than
$\left|\mathcal{V}^0(\boldsymbol{A},\boldsymbol{d})\right|\cdot\left|\mathcal{V}^r(\boldsymbol{A},\boldsymbol{d})\right|$.
Because $\left|\mathcal{V}^0(\boldsymbol{A},\boldsymbol{d})\right|\leq{m+n+1\choose n}$ and
$\displaystyle{\left|\mathcal{V}^r(\boldsymbol{A},\boldsymbol{d})\right|\leq {\left|\mathcal{V}^0(\boldsymbol{A},\boldsymbol{d})\right|\choose r+1}}$
according to Proposition \ref{prop:cover}, the number of executions of the third \texttt{for}-loop is polynomial in $m$.
% Further, $\left|\mathcal{V}^{r+1}\right|$ never exceeds $\left|\mathcal{V}^{r+1}(\boldsymbol{A},\boldsymbol{d})\right|$ (Lemma \ref{lem:alg}),
% implying that the search for $W\in\mathcal{V}^{r+1}$ also runs in polynomial time.
Finally, the dimensionality check of $P(V)$ is accomplished by searching $V$
for $r+2$ affinely independent vertices.
Thus, the number of arithmetic operations inside the loops is polynomial in $m$, which completes the proof.
\qed
\end{proof}


\subsection{A Polynomial Algorithm for the Pooling Problem with a Single Pool and Fixed Source Cardinality} \label{sec:p1pol}

By virtue of Algorithm \ref{alg:fpol} and Proposition \ref{prop:ptop}, it is provable that problem (\ref{eq:qobj})--(\ref{eq:tdom}) is solvable in polynomial time for bounded $|S|$.
Let $\boldsymbol{A}$ and $\boldsymbol{d}$ have one row for each $(t,k)\in T\times K$, and one column for each $s\in S_0$.
The entry in $\boldsymbol{d}$ corresponding to $(t,k)$ is set to $u_t^k$, and $a_{ij}$ is put equal to $q_s^k$,
where $i$ and $j$ are  the row and column corresponding to $(t,k)$ and $s$, respectively.
When receiving $\boldsymbol{Q}$ and $\boldsymbol{U}$ as input, Algorithm \ref{alg:fter}
outputs the set $\mathcal{T}(\mathcal{D},\boldsymbol{b})$ of sink sets that a feasible flow in $\mathcal{F}(\mathcal{D},\boldsymbol{b})$ possibly can satisfy.
% In all instances, we have $\mathcal{T}^+(\boldsymbol{0}^{|\mathcal{A}|})=T\in\mathcal{T}(\mathcal{D},\boldsymbol{b})$.
% However, the set $\mathcal{T}$ produced by Algorithm \ref{alg:fter} does not contain $T$ unless there exists some
% $\boldsymbol{f}\neq\boldsymbol{0}^{|\mathcal{A}|}$ for which $\mathcal{T}^+(\boldsymbol{f})=T$.

\begin{algorithm} \caption{\texttt{findTerminals}($\boldsymbol{Q}$, $\boldsymbol{U}$)} \label{alg:fter}
\begin{algorithmic} 
  \STATE Define the matrix $\boldsymbol{A}\in\mathbb{Q}^{(|T\times K|)\times|S_0|}$ and the vector $\boldsymbol{d}\in\mathbb{Q}^{|T\times K|}$
         such that they correspond to inequalities (\ref{eq:t}). Let $n\leftarrow |S_0|$.
  \STATE Find $\mathcal{V}^0,\ldots,\mathcal{V}^n$ by calling \texttt{findPolytopes}($\boldsymbol{A}$, $\boldsymbol{d}$)
  \STATE $\mathcal{T}\leftarrow\{T\}$ \COMMENT{all sinks are satisfied by $\boldsymbol{0}^{\left|\mathcal{A}\right|}$}
  \FOR[loop through all polytope dimensions $r$]{$r\leftarrow 0,\ldots,n$}
    \FOR[loop through all $r$-polytopes]{all $V\in\mathcal{V}^r$}
      \STATE $\displaystyle{\boldsymbol{x}\leftarrow\frac{1}{|V|}\sum_{\boldsymbol{y}\in V}\boldsymbol{y}}$ ~~\COMMENT{let $\boldsymbol{x}$ be any relative interior point of $P(V)$}
      \STATE $T^+\leftarrow T$ \COMMENT{assume $\boldsymbol{x}$ satisfies all concentration bounds of all sinks}
      \FOR[check component $k$]{all $k\in K$}
        \STATE $\displaystyle{T^+\leftarrow \left\{t\in T^+: \sum_{s\in S_0}q_s^kx_s \leq u_t^k\right\}}$ \COMMENT{keep only sinks that $\boldsymbol{x}$ satisfies}
      \ENDFOR
      \STATE $\mathcal{T}\leftarrow\mathcal{T}\cup\{T^+\}$ ~~\COMMENT{include the feasible sink set}
    \ENDFOR
  \ENDFOR
  \RETURN{$\mathcal{T}$}
\end{algorithmic}
\end{algorithm}

\begin{theorem} \label{prop:p1smax}
  For any $s_{\max}\geq 1$, the instance class of the \textsc{Pooling Problem} with $|I|=1$ and $|S|\leq s_{\max}$ is solvable in polynomial time.
\end{theorem}
\begin{proof}
According to Proposition \ref{prop:z}, the optimal cost in all instances in question is $\min\left\{z\left(T^+\right):T^+\in\mathcal{T}\right\}$,
where $\mathcal{T}$ is the output from Algorithm \ref{alg:fter}.
Because $z\left(T^+\right)$ is the solution to the compact LP (\ref{eq:qobj})--(\ref{eq:t-}),
the claim then follows directly from Proposition \ref{prop:ptop} with $n=|S|-1<s_{\max}$.
\qed
\end{proof}

\begin{corollary} \label{cor:ext}
The extension of Problem \ref{prob:pp} where a positive lower bound is imposed on the flow entering each sink in $T'\subseteq T$,
is solvable in polynomial time for instances where $|I|=1$ and $|S|\leq s_{\max}$.
\end{corollary}

\begin{proof}
Follows immediately by modifying Algorithm \ref{alg:fter} such that all sink sets $T^+\in\mathcal{T}(\mathcal{D},\boldsymbol{b})$ that do not contain $T'$ are discarded.
\qed
\end{proof}

It is also straightforward to adapt the proofs discussed in Sections \ref{sec:fewt}--\ref{sec:fewk} to the extension considered in Corollary \ref{cor:ext}.
In conclusion, the extended problem is solvable in polynomial time for the single-pool instance classes where either $|T|$, $|K|$ or $|S|$ is bounded by a constant.

Recall from Section \ref{sec:problem} that $\mathcal{D}$ is defined to have arcs only from $S$ to $I$ and from $I$ to $T$.
Given this restriction, instances that accept direct flow from a source $s$ to a sink $t$ contain a pool $i$ where $S_i=\{s\}$ and $T_i=\{t\}$.
Alternatively, $\mathcal{D}$ could be defined such that the direct arc from $s$ to $t$ is allowed, making the pool $i$ redundant.
Most studies \cite{haverly78,sahinidis05,dey15,alfaki13a} follow this convention.
It is however important to note that the results in Sections \ref{sec:onepool}--\ref{sec:i1sbnd}
are not proved to hold if $|I|=1$ in this alternative definition of $\mathcal{D}$.
That is, it is an open question whether the \textsc{Pooling Problem} is solvable in polynomial time when either $|T|$, $|K|$ or $|S|$ is fixed,
and only one pool has more than one entering arc and more than one leaving arc.

\begin{remark}
After submission of the current text, Boland et al.\ \cite{boland15} submitted a manuscript with an alternative proof of Proposition \ref{prop:p1smax}.
Their analysis has points in common with the analysis in Section \ref{sec:simp},
as it also considers partitions of the unit simplex induced by a set of linear inequalities.
Moreover, \cite{boland15} proves that in instances where $|I|=1$, the \textsc{Pooling Problem} is solvable in terms of at most
$\sum_{j=0}^{|S|-1}{|K|\choose j}|T|^j$ compact linear programs. \qed
\end{remark}

%-------------------------------------------------
\section{The Pooling Problem with a Single Quality Parameter} \label{sec:k1}

In the remainder of the text, we study a rich instance class where $|I|$ is arbitrary, $|K|=1$, and $|S|=|T|=2$.
By a reduction from \textsc{Exact Cover By 3-sets}, it is proved \cite{haugland15a} that the \textsc{Pooling Problem} remains strongly NP-hard when $|K|=1$.
When also $|S|=|T|=2$, at least weak NP-hardness persists.
However, the reduction \cite{haugland15a} from \textsc{Packing In Two Bins} for the $|S|=|T|=2$ case, does not prove strong NP-hardness.
\textsc{Packing In Two Bins} asks whether a set of items of given integer weights fit into two bins with identical integer capacities.
There exists a dynamic programming algorithm, which runs in polynomial time for instance classes where item weights and
bin capacities are bounded by some polynomial function of the input size.
In this section, we show how the dynamic programming algorithm for \textsc{Packing In Two Bins} can be adapted to
the instance class of the \textsc{Pooling Problem} under study.

\subsection{A Dynamic Programming Algorithm} \label{sec:dp}
Each recursive step of the dynamic programming algorithm for \textsc{Packing In Two Bins}
amounts to selecting the bin into which a given item should be put.
The corresponding operation in the algorithm given below, is to select in which out of five possible subsets a given pool belongs.
Each of four of the subsets corresponds to a particular source or sink,
and all pools allocated to the set receive (send) zero flow from (to) the source (sink) in question.
Pools allocated to the fifth set receive (send) positive flow from (to) both sources (sinks).

It is shown below that the collection of pools in one of the former four sets can be replaced by a new pool with capacity equal to the total capacity in the set.
Further, it is shown that under the assumptions made, it is optimal to have at most one pool in the latter set.
Hence, the dynamic programming algorithm reduces the network to an instance with only five pools.
Instances of this small, fixed size are shown to have closed-form solutions, and they are thereby solvable in constant time.
Since a running time analysis of the dynamic programming algorithm proves that the reduction runs in pseudo-polynomial time,
we thus have an algorithm with pseudo-polynomial running time for $|K|=1$ and $|S|=|T|=2$.

Because $|K|=1$, indices corresponding to components are omitted in all notation used.
Also for the purpose of simplified notation, we associate the nodes $\mathcal{N}$ with integers,
and for reasons to become clear later, we let $S=\{0,1\}$ and $T=\{7,8\}$.
Without loss of generality \cite[Proposition 2]{haugland15a}, we assume $q_0=0$, $q_1=1$, and $u_7\leq u_8$.
Further, we focus on instances with pool independent costs, i.e. all arcs leaving (entering) the same source (sink) have identical costs.
We denote by $c_s$ and $c_t$ the unit costs of arcs
$(s,i)\in \mathcal{A}_S$ and $(i,t)\in \mathcal{A}_T$, respectively.
For any feasible flow, we define the \emph{flow degree} of a node as the number of incident arcs carrying positive flow.

\begin{proposition} \label{prop:k1}
\textsc{Pooling Problem} instances with $|K|=1$, $|S|=|T|=2$, and pool independent costs have optimal solutions where at most one pool has flow degree 4.
\end{proposition}

\begin{proof}
Let $\boldsymbol{f}$ be an optimal flow, and assume $f_{si},f_{sj},f_{it},f_{jt}>0$ for all $s\in S$ and $t\in T$, where $i,j\in I$ and $i\neq j$.
Consider the cycles $C_1=(1,j,0,i,1)$ and $C_2=(j,8,i,7,j)$ consisting of forward and backward arcs in $\mathcal{D}$.
Flow augmentations $\varepsilon$ and $\delta$ along cycles $C_1$ and $C_2$, respectively,
yield a flow denoted $\boldsymbol{g}$, where $\boldsymbol{g}=\boldsymbol{f}$ if $\varepsilon=\delta=0$
(the dependency of $\varepsilon$ and $\delta$ in $\boldsymbol{g}$ is suppressed for the purpose of simplified notation).
Let $\Omega$ be the set of pairs $(\varepsilon,\delta)$ for which $\boldsymbol{g}\geq \boldsymbol{0}^{|\mathcal{A}|}$.
That is, $\Omega$ consists of all $(\varepsilon,\delta)$ satisfying
\begin{eqnarray*}
  & -\min\left\{f_{0i},f_{1j}\right\}\leq\varepsilon\leq\min\left\{f_{0j},f_{1i}\right\}, & \mbox{and} \\
  & -\min\left\{f_{i7},f_{j8}\right\}\leq\delta\leq\min\left\{f_{j7},f_{i8}\right\}.
\end{eqnarray*}
Because the augmentations do not affect the capacity consumption at the nodes, $\boldsymbol{g}\in\mathcal{F}(\mathcal{D},\boldsymbol{b})$ if $(\varepsilon,\delta)\in\Omega$.
Since the costs are pool independent, also the objective function values at $\boldsymbol{f}$ and $\boldsymbol{g}$ are identical.
It follows that $\boldsymbol{g}$ is optimal if it satisfies the concentration bounds.

Let $\boldsymbol{w}\in\mathbb{R}^{|\mathcal{N}|}$ and $\boldsymbol{v}\in\mathbb{R}^{|\mathcal{N}|}$
be the concentration vectors corresponding to $\boldsymbol{f}$ and $\boldsymbol{g}$, respectively.
We prove the proposition by showing that there exists an augmentation $(\varepsilon,\delta)\in\Omega$ such that $\boldsymbol{g}$
is zero on one arc of $C_1$ or $C_2$, and such that $v_{7}=w_7$ and $v_{8}=w_8$.

Let $F_{\ell}=f_{\ell 7}+f_{\ell 8}=g_{\ell 7}+g_{\ell 8}$ be the total flow through pool $\ell \in I$.
From $q_0=0$ and $q_1=1$, it follows that $w_i=f_{1i}/F_i$ and $w_j=f_{1j}/F_j$. We also have
\[
  \mbox{$\displaystyle{v_i = \frac{w_iF_i-\varepsilon}{F_i}=\frac{f_{1i}-\varepsilon}{F_i}}$, and
  $\displaystyle{v_j = \frac{w_jF_j+\varepsilon}{F_j}=\frac{f_{1j}+\varepsilon}{F_j}}$, yielding}
\]
\[
  \displaystyle{v_7 = \frac{w_7F_7-w_if_{i7}-w_jf_{j7}+v_ig_{i7}+v_jg_{j7}}{F_7} =}
\]
\[
      \displaystyle{w_7 +\frac{-f_{1i}f_{i7}/F_i - f_{1j}f_{j7}/F_j + (f_{1i}-\varepsilon)(f_{i7}+\delta)/F_i + (f_{1j}+\varepsilon)(f_{j7}-\delta)/F_j}{F_7}}.
\]

\noindent
It follows that $v_7=w_7$ if
\begin{equation}
  \varepsilon\delta\left(F_i+F_j\right) + \varepsilon\left(F_jf_{i7}-F_if_{j7}\right) + \delta\left(F_if_{1j}-F_jf_{1i}\right) = 0. \label{eq:propk1}
\end{equation}

The set of pairs $(\varepsilon,\delta)$ satisfying (\ref{eq:propk1}) is the union of two continuous curves $H_1$ and $H_2$.
If $F_jf_{i7}-F_if_{j7}\neq 0\neq F_if_{1j}-F_jf_{1i}$, then $H_1\cup H_2$ is a hyperbola.
Otherwise, $H_1\cup H_2$ consists of two straight lines, one of which is a coordinate axis (where either $\varepsilon=0$ or $\delta=0$).
Because neither $H_1\setminus\Omega$ nor $H_2\setminus\Omega$ is empty, and $(0,0)\in \left(H_1\cup H_2\right)\cap\Omega$,
continuity implies that $H_1\cup H_2$ intersects the boundary of $\Omega$.
It follows that there exists some $(\varepsilon,\delta)$ such that $\boldsymbol{g}$ is zero on one of the arcs of $C_1$ or $C_2$,
and such that $v_7=w_7$.
Because
$\sum_{\ell\in I_s}f_{s\ell}=\sum_{\ell\in I_s}g_{s\ell}$ ($s\in S$) and
$\sum_{\ell\in I_t}f_{\ell t}=\sum_{\ell\in I_t}g_{\ell t}$ ($t\in T$), (\ref{eq:totqual}) implies that also
$v_8=w_8$,
which completes the proof. \qed
\end{proof}

% \begin{centering}
\begin{figure}[ht]
\begin{center}
% \vspace{-3em}
\vspace{2em}
\includegraphics[scale=0.4]{base.pdf}
% \vspace{-18em}
\end{center}
\caption{Atomic instance}
\vspace{-2em}
\label{fig:atom}
\end{figure}
% \vspace{-2em}
% \end{centering}

Motivated by Proposition \ref{prop:k1}, we search for the optimal partition of the pools $I$ into 5 subsets denoted $J_2,\ldots,J_6$.
The subset $J_2$ ($J_6$)  consists of pools that do not receive flow from source 1 (source 0),
and $J_3$ ($J_5$) consists of pools the do not send flow to sink 8 (sink 7).
Subset $J_4$ is either empty, or consists of the unique pool of flow degree 4.
Ties are broken arbitrarily.
Associated with the partition $(J_2,\ldots,J_6)$, we define another \textsc{Pooling Problem} instance with directed graph
$\bar{\mathcal{D}}=\left(\bar{\mathcal{N}}, \bar{\mathcal{A}}\right)$
(see Fig.\ \ref{fig:atom}), henceforth referred to as an \emph{atomic} instance.
The node set $\bar{\mathcal{N}}=\bar{S}\cup\bar{I}\cup\bar{T}$ consists of the sources $\bar{S}=S=\{0,1\}$,
pools $\bar{I}=\{2,\ldots,6\}$, and sinks $\bar{T}=T=\{7,8\}$.

Each pool $\ell\in\bar{I}$ corresponds to one of the sets $J_2,\ldots,J_6$.
The arc set
$\bar{\mathcal{A}}$ is hence defined such that the neighbors of sources 0 and 1 and sinks 7 and 8 are
$\bar{I}_0=\bar{I}\setminus\{6\}$,
$\bar{I}_1=\bar{I}\setminus\{2\}$,
$\bar{I}_7=\bar{I}\setminus\{5\}$, and
$\bar{I}_8=\bar{I}\setminus\{3\}$, respectively.
Costs $\boldsymbol{c}$, source concentrations $\boldsymbol{q}$, and sink concentration bounds $\boldsymbol{u}$ are identical to those of the original instance.
The same is true for source and sink capacities, whereas the capacity of pool $\ell\in\bar{I}$
is the cumulative capacity $B_{\ell}=\sum_{i\in J_{\ell}}b_i$ of pools in $J_{\ell}$.
Denote the node capacity vector by $\boldsymbol{B}$, where $B_s=b_s$ and $B_t=b_t$ ($s\in \bar{S}$, $t\in \bar{T}$).
For the case of notational brevity, the dependency of $(J_2,\ldots,J_6)$ in $\boldsymbol{B}$ is suppressed.

Let $\zeta\left(\mathcal{D},\boldsymbol{b}\right)$ denote the minimum cost in the given instance.
Relevance of the atomic instances is justified as follows:

\begin{proposition} \label{prop:atom}
\textsc{Pooling Problem} instances with $|K|=1$, $|S|=|T|=2$ and pool independent costs,
have $\zeta\left(\mathcal{D},\boldsymbol{b}\right)=\min\zeta\left(\bar{\mathcal{D}},\boldsymbol{B}\right)$,
where the minimum is defined over all partitions $(J_2,\ldots,J_6)$ of $I$ such that $|J_4|\leq 1$.
\end{proposition}

\begin{proof}
To prove $\zeta\left(\mathcal{D},\boldsymbol{b}\right)\geq\min\zeta\left(\bar{\mathcal{D}},\boldsymbol{B}\right)$,
let $\boldsymbol{f}$ be an optimal flow in $\mathcal{D}$ with corresponding partition $(J_2,\ldots,J_6)$ of $I$.
Without loss of generality (Proposition \ref{prop:k1}), assume $|J_4|\leq 1$.
Let $g_{s\ell}=\sum_{i\in J_{\ell}}f_{si}$ ($s\in \bar{S},\ell\in\bar{I}_s$) and \\
$g_{\ell t}=\sum_{i\in J_{\ell}}f_{it}$ ($t\in \bar{T},\ell\in\bar{I}_t$).
Because $B_{\ell}=\sum_{i\in J_{\ell}}b_i$ and $\boldsymbol{f}\in\mathcal{F}\left(\mathcal{D},\boldsymbol{b}\right)$,
it follows immediately that $g\in\mathcal{F}\left(\bar{\mathcal{D}},\boldsymbol{B}\right)$ and
\[\sum_{s\in\bar{S}}c_s\sum_{\ell\in\bar{I}_s}g_{s\ell}+\sum_{t\in\bar{T}}c_t\sum_{\ell\in\bar{I}_t}g_{\ell t}
=\sum_{(s,i)\in\mathcal{A}_S}c_sf_{si}+\sum_{(i,t)\in\mathcal{A}_T}c_tf_{it}.\]

It only remains to prove that $\boldsymbol{g}$ satisfies the concentration constraints.
Let $\boldsymbol{w}$ and $\boldsymbol{v}$ be concentration vectors corresponding to $\boldsymbol{f}$ and $\boldsymbol{g}$, respectively.
Thus, $v_2=w_i=0$ ($i\in J_2$) and $v_6=w_i=1$ ($i\in J_6$).
Pool concentration $v_{\ell}$ is a convex combination of $\left\{w_i\right\}_{i\in J_{\ell}}$ ($\ell=3,5$), given by
$v_3=\sum_{i\in J_3}w_i\frac{f_{i7}}{g_{37}}$ and
$v_5=\sum_{i\in J_5}w_i\frac{f_{i8}}{g_{58}}$.
With respect to the concentration $v_4$, if $|J_4|=1$, we have $g_{s4}=f_{si'}$ for $i'\in J_4$ and $s\in \bar{S}$, yielding $v_4=w_{i'}$.
Otherwise, $v_4$ is arbitrary.
Then $\displaystyle{\sum_{\ell\in\bar{I}_7}v_{\ell}g_{\ell 7}}$ equals
\[
  \displaystyle{v_3g_{37}+v_4g_{47}+v_6g_{67}
  = \sum_{i\in J_3}w_i\frac{f_{i7}}{g_{37}}g_{37} + w_{i'}f_{i'7} + \sum_{i\in J_6}w_if_{i7} = \sum_{i\in I_7}w_if_{i7}}.
\]

\noindent
It follows that
\[
  \displaystyle{v_7 = \frac{\sum_{\ell\in\bar{I}_7}v_{\ell}g_{\ell 7}}{\sum_{\ell\in\bar{I}_7}g_{\ell 7}} = \frac{\sum_{i\in I_7}w_if_{i7}}{\sum_{i\in I_7}f_{i7}}=w_7\leq u_7},
\]
\noindent
since $\boldsymbol{f}$ is feasible.
Fulfillment of $v_8\leq u_8$ is proved analogously.

To prove $\zeta\left(\mathcal{D},\boldsymbol{b}\right)\leq\min\zeta\left(\bar{\mathcal{D}},\boldsymbol{B}\right)$,
let $\boldsymbol{g}$ be an optimal flow in $\bar{\mathcal{D}}$.
For \\ $(s,\ell)\in\{(0,2),(1,6)\}$, let $f_{si}=g_{s\ell}\cdot b_i/B_{\ell}$ and $f_{it}=f_{si}\cdot g_{\ell t}/g_{s\ell}$ ($i\in J_{\ell}$, $t\in T$).
For $(\ell,t)\in\{(3,7),(5,8)\}$, let $f_{it}=g_{\ell t}\cdot b_i/B_{\ell}$ and $f_{si}=f_{it}\cdot g_{s\ell}/g_{\ell t}$ ($i\in J_{\ell}$, $s\in S$).
If $|J_4|=1$, let $f_{si'}=g_{s4}$ and $f_{i't}=g_{4t}$, ($s\in S$, $t\in T$, $i'\in J_4$).
Utilizing $B_{\ell}=\sum_{i\in J_{\ell}}b_i$ and feasibility of $\boldsymbol{g}$, the proof of feasibility of $\boldsymbol{f}$ is analogous to the first part of the proof.
Because the objective function values of $\boldsymbol{f}$ and $\boldsymbol{g}$ in their respective problems are equal, the proof is complete.
\qed
\end{proof}

Observe that $|J_4|\leq 1$ is an essential condition in Proposition \ref{prop:atom}.
If rather $|J_4|>1$, letting $g_{s4}=\sum_{i\in J_4}f_{si}$ and $g_{4t}=\sum_{i\in J_4}f_{it}$ ($s\in S$, $t\in T$) does not necessarily yield $v_t=w_t$.
Aggregating pools in $J_{\ell}$ into one common pool is valid only if either their in-degree or out-degree is one.
This is closely related to the fact that the \textsc{Pooling Problem} is an LP in case $\min\left\{|S_i|,|T_i|\right\}=1$ for all pools $i\in I$ \cite[Proposition 3]{haugland15a}.

We now go on to derive a dynamic programming algorithm for reducing the original instance to an atomic instance corresponding to an optimal partition $(J_2,\ldots,J_6)$.
For simplified exposition, we assume henceforth that all pool capacities $b_i$ are integers.
Denote the original pools by $i_1,\ldots,i_m$, where $m=|I|$.
At stage $n=1,\ldots,m,m+1$ of the dynamic program, each of the pools $i_{1},\ldots,i_{n-1}$ has been assigned to some set $J_{\ell}$,
whereas pools $i_{n},\ldots,i_m$ are left to be assigned.
The state variable is the capacity vector $\boldsymbol{B}$, where $B_{\ell}$ is the total capacity currently assigned to $J_{\ell}$.

The cost function $z_n(\boldsymbol{B})$ represents the minimum obtainable cost given state $\boldsymbol{B}$ at stage $n$.
% That is, $z_n(D)$ can be attained if $i_{1},\ldots,i_{n-1}$ are assigned to one of $I_2,\ldots,I_6$
% such that pools of total capacity $B_{\ell}$ are assigned to $I_{\ell}$ (${\ell}=2,\ldots,6$).
At stage $m+1$, all pools are allocated, and $z_{m+1}(\boldsymbol{B})$ is the minimum cost of an atomic instance with pool capacities $\boldsymbol{B}$.
The assignment of $i_{n}$ made at stage $n\leq m$, is represented by a unit vector in the feasible set
\[
  E(\boldsymbol{B})=\left\{\boldsymbol{e}=(e_2,\ldots,e_6)\in\{0,1\}^5, \sum_{{\ell}=2}^6e_{\ell}=1, \mbox{$e_4=0$ if $B_4>0$}\right\}.
\]
Optimizing the assignment yields the backward recursion formula
$$
  z_n(\boldsymbol{B})=\min_{\boldsymbol{e}\in E(\boldsymbol{B})} z_{n+1}(\boldsymbol{B}+b_n\boldsymbol{e}),
$$
\noindent
where $n=1,\ldots,m$, and $z_1(0)$ equals the optimal objective function value $\zeta(\mathcal{D},\boldsymbol{b})$.
Correspondingly, the dynamic programming procedure depicted in Algorithm \ref{alg:dp} applies.

\begin{algorithm} \caption{\texttt{DP}($S$,$I$,$T$,$\mathcal{A}$,$\boldsymbol{c}$,$\boldsymbol{b}$,$\boldsymbol{q}$,$\boldsymbol{u}$)} \label{alg:dp}
\begin{algorithmic} 
  \FOR{$\forall \boldsymbol{B}\in\mathbb{N}^5:\sum_{{\ell}=2}^6B_{\ell}=\sum_{i\in I}b_i$}
    \STATE $z_{m+1}(\boldsymbol{B})\leftarrow$ minimum cost in the atomic instance with pool capacities $\boldsymbol{B}$
  \ENDFOR
  \FOR{$n\leftarrow m,\ldots,1$}
    \FOR{$\forall \boldsymbol{B}\in\mathbb{N}^5:\sum_{{\ell}=2}^6B_{\ell}=\sum_{j=1}^{n-1}b_{i_j}$}
      \STATE $z_n(\boldsymbol{B})=\min\left\{z_{n+1}(\boldsymbol{B}+b_n\boldsymbol{e}):\boldsymbol{e}\in E(\boldsymbol{B})\right\}$
     \ENDFOR
  \ENDFOR
  \RETURN{$z_1(0)$}
\end{algorithmic} 
\end{algorithm}

\subsection{Running Time Analysis} \label{sec:dptime}
In the first loop of Algorithm \ref{alg:dp}, $N={4+\sum_{i\in I}b_i\choose 4}$ atomic instances are solved.
The remainder of the algorithm runs in $\mathcal{O}(N|I|)$ time.
Despite its fixed size, it is not obvious that an atomic instance is solvable by a constant number of arithmetic operations.
Making a reasonable assumption about the arc costs is however a sufficient condition:

\begin{assumption}[Quality consistent costs] The arc costs $\boldsymbol{c}$ satisfy: \label{as:cost}
\begin{enumerate}
  \item $c_0>c_1$, \label{as:s}
  \item $c_8>c_7$, \label{as:t}
  \item $c_0+c_8>0$, and \label{as:st}
  \item $(1-u_t)c_0+u_tc_1+c_t<0$, $t=7,8$. \label{as:eq}
\qed
\end{enumerate}
\end{assumption}

\noindent
Assumption \ref{as:cost} is likely to hold in many practical instances.
The first two points state, respectively, that the source supplying the better concentration is also the more expensive,
and that the sales prices are higher at the sink with the stricter concentration requirement.
Further, Assumption \ref{as:cost}.\ref{as:st} states that it is not profitable to produce the poorer end product
(at sink 8) uniquely by use of the more expensive supply (source 0).
Finally, the last point states that at both sinks, a blend satisfying the concentration constraints by equality is profitable.

For atomic instances fulfilling Assumption \ref{as:cost}, a solution algorithm with proven convergence in
a constant number of operations on rational numbers has recently been given \cite{haugland15c}.
Assumptions \ref{as:cost}.\ref{as:s}--\ref{as:cost}.\ref{as:st} are made merely in order to focus the study,
and to limit the number of cases to be investigated.
If Assumption \ref{as:cost}.\ref{as:eq} does not hold, at least one of the sinks can be removed, and the resulting \textsc{Pooling Problem} instance becomes a linear program.
Quality consistent instances do not appear to be easier than others, and we therefore conjecture that the result in \cite{haugland15c} also applies to
instances that do not fulfill Assumption \ref{as:cost}.
The proof of \cite[Theorem 3]{haugland15a} shows that the \textsc{Pooling Problem} with $|K|=1$ and $|S|=|T|=2$ is NP-hard even when the costs are
pool independent and concentration consistent.

Let $\nu\left(\mathcal{D},\boldsymbol{b},\boldsymbol{c},\boldsymbol{q},\boldsymbol{u}\right)$ denote the size of a problem instance.
It follows from the running time analysis that for any polynomial $p:\mathbb{Z}\mapsto\mathbb{Z}$,
there exists a solution algorithm with polynomial running time for the class of \textsc{Pooling Problem} instances
where $|K|=1$, $|S|=|T|=2$, costs are pool independent and concentration consistent, and $\sum_{i\in I}b_i\leq p\left(\nu\left(\mathcal{D},\boldsymbol{b},\boldsymbol{c},\boldsymbol{q},\boldsymbol{u}\right)\right)$.
We have thereby proved:
\begin{theorem} \label{prop:pseudo}
The \textsc{Pooling Problem} with $|K|=1$, $|S|=|T|=2$, and pool independent costs satisfying Assumption \ref{as:cost} is solvable in pseudo-polynomial time.
\end{theorem}


%---------------------------------------------------
\section{Conclusions}\label{sec:concl}
This work presents polynomial running time algorithms for two instance classes of the \textsc{Pooling Problem}.
In instances where the flow network contains only one pool, and an upper bound on the source cardinality applies,
the problem can be solved in polynomial time by solving an associated partitioning problem:
Hyperplanes corresponding to concentration constraints intersect the unit simplex in the space of the source proportion variables.
A list of the polytopes into which the simplex is partitioned, yields a list of linear programs that can be solved in order to solve the original problem.
Polynomial running time of this algorithm is formally proved by proving that the length of the list is polynomial in the number of constraints.

Instances with only one component and two sources and sinks are solved by a dynamic programming algorithm.
This algorithm is based upon a reduction of general instances to instances of fixed size.
When conditions that are likely to apply in practice are assumed,
it is proved that the instances of constant size are solvable by a constant number of arithmetic operations.
From this it is directly concluded that the dynamic programming algorithm has pseudo-polynomial running time.

The contribution offered by the two algorithms under study is mainly of a theoretical character.
Questions on the computational complexity of the \textsc{Pooling Problem} that in earlier studies were left open, are now answered.
In their present form, the algorithms appear to be computationally demanding, despite their favorable theoretical properties.
How, and to what extent, the algorithmic ideas presented in this work can be made suitable for solving practical instances,
is left as a challenge for future research.

Applications of the algorithms can possibly be found in the design of heuristic methods for more general instance classes than those considered in the present text.
Algorithms tailored for the one-pool case can be applied in order to construct feasible flows in networks with multiple pools.
This follows from the obvious fact that any feasible flow using only one pool is feasible even if other pools are present in the network.
Further, any convex combination of feasible flows using distinct pools is also feasible.
Taking advantage of these observations, construction heuristics for multi-pool instances can be designed.

Heuristics for improving a given feasible solution can also be designed by utilizing algorithmic elements presented in this paper.
The idea of flow augmenting paths \cite{haugland15c} and cycles can be generalized to instances with more than one component.
Improvement heuristics for general instances can thus be based upon
identification of subnetworks in which the flow can be augmented, such that costs are reduced and feasibility is preserved.
Exploiting these ideas in order to find computationally efficient solution methods is left as a challenge for future research.

\subsection*{Acknowledgment}
\small{The work of the second author has been funded by grants from the Spanish state (TIN2012-37483-C03-01 and TIN2015-66680), Junta de Andaluc\'{\i}a (P11-TIC-7176),
in part financed by the European Regional Development Fund (ERDF).}


%\bibliographystyle{abbrvnat}
%\bibliography{duopbib}
\begin{thebibliography}{99}
\bibitem{stigler45}
Stigler, G.J.:
\newblock The cost of subsistence.
\newblock J. Farm Econ.
\newblock \textbf{27}(2), 303--314 (1945)

\bibitem{dantzig90}
Dantzig, G.B.:
\newblock The diet problem.
\newblock Interfaces
\newblock \textbf{20}(4), 43--47 (1990)

\bibitem{haverly78}
Haverly, C.A.:
\newblock Studies of the behavior of recursion for the pooling problem.
\newblock ACM SIGMAP Bulletin,
\newblock \textbf{25}, 19--28 (1978)

\bibitem{baker85}
Baker, T.E., Lasdon, L.S.:
\newblock Successive linear programming at Exxon.
\newblock Manage.\ Sci.
\newblock \textbf{31}(3), 264--274 (1985)

\bibitem{dewitt89}
DeWitt, C.W., Lasdon, L.S., Waren, A.D., Brenner, D.A., Melham, S.:
\newblock OMEGA: An improved gasoline blending system for Texaco.
\newblock Interfaces
\newblock \textbf{19}(1), 85--101 (1989)

\bibitem{floudas90b}
Floudas, C.A., Visweswaran, V.:
\newblock A global optimization algorithm (GOP) for certain classes of nonconvex NLPs: I. Theory.
\newblock Comput.\ Chem.\ Eng.
\newblock \textbf{14}(12), 1397--1417 (1990)

\bibitem{audet04}
Audet, C., Brimberg, J., Hansen, P., Le Digabel, S., Mladenovi\'{c}, N.:
\newblock Pooling problem: alternate formulations and solution methods.
\newblock Manage.\ Sci.
\newblock \textbf{50}(6), 761--776 (2004)

\bibitem{foulds92}
Foulds, L.R., Haugland, D., J\"{o}rnsten, K.:
\newblock A bilinear approach to the pooling problem.
\newblock Optimization,
\newblock \textbf{24}, 165--180 (1992)

\bibitem{misener11}
Misener, R., Thompson, J.P, Floudas, C.A.:
\newblock APOGEE: Global optimization of standard, generalized, and extended pooling
problems via linear and logarithmic partitioning schemes.
\newblock Comput.\ Chem.\ Eng.
\newblock \textbf{35}, 876--892 (2011)

\bibitem{sahinidis05}
Sahinidis, N.V., Tawarmalani, M.:
\newblock Accelerating branch-and-bound through a modeling language construct for relaxation-specific constraints.
\newblock J.\ Global Optim.
\newblock \textbf{32}(2), 259--280 (2005)

\bibitem{visweswaran96b}
Visweswaran, V., Floudas, C.A.:
\newblock Computational results for an efficient implementation of the GOP algorithm and its variants.
\newblock In: Grossmann, I.E.\ (ed.): Global Optimization in Chemical Engineering, pp.\ 111--153.
\newblock Kluwer Academic Publishers, Dordrecht, The Netherlands (1996)

\bibitem{adhya99}
Adhya, N., Tawarmalani, M., Sahinidis, N.V.:
\newblock A Lagrangian approach to the pooling problem.
\newblock Ind.\ Eng.\ Chem.\ Res.
\newblock \textbf{38}(5), 1965--1972 (1999)

\bibitem{almutairi09}
Almutairi, H., Elhedhli, S.:
\newblock A new Lagrangian approach to the pooling problem.
\newblock J.\ Global Optim.
\newblock \textbf{45}(2), 237--257 (2009)

\bibitem{bental94}
Ben-Tal, A., Eiger, G., Gershovitz, V.:
\newblock Global minimization by reducing the duality gap.
\newblock Math.\ Program.
\newblock \textbf{63}(2), 193--212 (1994)

\bibitem{dey15}
Dey, S.S., Gupte, A.:
\newblock Analysis of MILP techniques for the pooling problem.
\newblock Oper.\ Res., \textbf{62}(2), 412--427 (2015)

\bibitem{kallrath05}
Kallrath, J.:
\newblock Solving planning and design problems in the process industry using mixed integer and global optimization.
\newblock Ann.\ Oper.\ Res.
\newblock \textbf{140}(1), 339--373 (2005)

\bibitem{galan98}
Galan, B., Grossmann, I.E.:
\newblock Optimal design of distributed wastewater treatment networks.
\newblock Ind.\ Eng.\ Chem.\ Res.
\newblock \textbf{37}(10), 4036--4048 (1998)

\bibitem{misener10b}
Misener, R., Floudas, C.A.:
\newblock Global optimization of large-scale generalized pooling problems: Quadratically constrained MINLP models.
\newblock Ind.\ Eng.\ Chem.\ Res.
\newblock \textbf{49}, 5424--5438 (2010)

\bibitem{misener09b}
Misener, R., Floudas, C.A.:
\newblock Advances for the pooling problem: modeling, global optimization, and computational studies.
\newblock Appl.\ Comput.\ Math.
\newblock \textbf{8}, 3--22 (2009)

\bibitem{alfaki13a}
Alfaki, M., Haugland, D.:
\newblock Strong formulations for the pooling problem.
\newblock J.\ Global Optim.
\newblock \textbf{56}(3), 897--916 (2013)

\bibitem{haugland15a}
Haugland, D.:
\newblock The computational complexity of the pooling problem.
\newblock To appear in J.\ Global Optim.
\newblock DOI 10.1007/s10898-015-0335-y (2015)

\bibitem{aleksandrov56}
Aleksandrov, P.S.:
\newblock Combinatorial Topology, Vol.\ 1.
\newblock Graylock Press, Rochester, New York (1956)

\bibitem{deloera10}
De Loera, J.A., Rambau, J., Santos, F.:
\newblock Triangulations - Structures for Algorithms and Applications.
\newblock Springer Verlag, Berlin, Germany (2010)

\bibitem{edmonds70}
Edmonds, A.L.:
\newblock Simplicial decompositions of convex polytopes.
\newblock Pi Mu Epsilon Journal,
\newblock \textbf{5}(3), 124--128 (1970)

\bibitem{murty95}
Murty, K.G., Chung, S.-J.:
\newblock Segments in enumerating faces.
\newblock Math.\ Program.
\newblock \textbf{70}(1), 27--45 (1995)

\bibitem{fukuda97}
Fukuda, K., Liebling, T.M., Margot, F.:
\newblock Analysis of backtrack algorithms for listing all vertices and all faces of a convex polyhedron.
\newblock Comp. Geom. Theor. Appl.
\newblock \textbf{8}, 1--12 (1997)

\bibitem{boland15}
Boland, N., Kalinowski, T., Rigterink, F.:
\newblock A polynomially solvable case of the pooling problem.
\newblock arXiv:1508.03181v2 [math.OC] (2015)

\bibitem{haugland15c}
Haugland, D., Hendrix, E.M.T.:
\newblock On a pooling problem with fixed network size.
\newblock In: Corman, F., Vo\ss, S., Negenborn,R.R.\ (eds.):
\newblock Computational Logistics, 6th International Conference, ICCL 2015, Proceedings,
\newblock Lect. Notes Comput. Sc. \textbf{9335}, 328--342 (2015)
\end{thebibliography}

\newpage
\section*{Appendix: Table of Notation}
\begin{description}
\item $\mathbb{R}^n$ = the $n$-dimensional space of real vectors
\item $\mathbb{Q}^n$ = the $n$-dimensional space of rational vectors
\item $\mathcal{D}=\left(\mathcal{N},\mathcal{A}\right)$ = directed graph with node set $\mathcal{N}$ and arc set $\mathcal{A}$
\item $S$ = set of source nodes
\item $I$ = set of pool nodes
\item $T$ = set of sink nodes
\item $K$ = set of components
\item $S_i$ = set of neighbor sources of pool $i$
\item $T_i$ = set of neighbor sinks of pool $i$
\item $I_j$ = set of neighbor pools of node $j$
\item $\mathcal{A}_S$ = set of source incident arcs
\item $\mathcal{A}_T$ = set of sink incident arcs
\item $\boldsymbol{Q}=\left(q_s^k\right)_{s\in S, k\in K}$ = matrix of source concentrations
\item $\boldsymbol{U}=\left(u_t^k\right)_{t\in T, k\in K}$ = matrix of sink concentration bounds
\item $\boldsymbol{b}$ = vector of node capacities
\item $\boldsymbol{c}$ = vector of arc costs
\item $s_0$ = a source for which $q_{s_0}^k=0$ for all $k\in K$
\item $S_0=S\setminus\{s_0\}$
\item $\boldsymbol{f}$ = flow vector
\item $\boldsymbol{W}=\left(w_j^k\right)_{j\in \mathcal{N}, k\in K}$ = matrix of flow concentrations
\item $\mathcal{F}\left(\mathcal{D},\boldsymbol{b}\right)$ = set of all flow vectors satisfying conservation of flow and node capacities
\item $\mathcal{T}^+:\mathcal{F}\left(\mathcal{D},\boldsymbol{b}\right)\Rightarrow 2^T$
      maps flow vector $\boldsymbol{f}$ to the set of sinks with all concentration constraints satisfied by $\boldsymbol{f}$
\item $\mathcal{T}\left(\mathcal{D},\boldsymbol{b}\right)$ = image of $\mathcal{T}^+$
\item $x_s$ = proportion of the flow through a pool originating from source $s$
\item $z\left(T^+\right)$ = minimum cost when only sinks in $T^+$ can receive positive flow
\item $P$ = polytope in $\mathbb{R}^n$
\item $V$ = finite set in $\mathbb{R}^n$
\item $V(P)$ = vertex set of polytope $P$
\item $P(V)$ = convex hull of $V$
\item $\boldsymbol{A}=\left(a_{ij}\right)_{i\in M, j \in N}$ = rational matrix
\item $\boldsymbol{d}=\left(d_i\right)_{i\in M}$ = rational vector
\item $\left(\bar{\boldsymbol{A}},\bar{\boldsymbol{d}}\right) = \left(\boldsymbol{A},\boldsymbol{d}\right)$ extended by rows with coefficients of the inequalities defining the unit simplex
\item $\left(M^<,M^>\right)$ = valid cover
\item $\Delta^n$ = unit $n$-simplex
\item $\mathcal{P}\left(\boldsymbol{A},\boldsymbol{d}\right)$ = polyhedral complex induced by $\left(\boldsymbol{A},\boldsymbol{d}\right)$
\item $\mathcal{P}^r\left(\boldsymbol{A},\boldsymbol{d}\right)$ = set of $r$-dimensional polyhedra in $\mathcal{P}\left(\boldsymbol{A},\boldsymbol{d}\right)$
\item $\mathcal{V}\left(\boldsymbol{A},\boldsymbol{d}\right)$ = vertex sets of polyhedra in $\mathcal{P}\left(\boldsymbol{A},\boldsymbol{d}\right)$
\item $\mathcal{V}^r\left(\boldsymbol{A},\boldsymbol{d}\right)$ = vertex sets of polyhedra in $\mathcal{P}^r\left(\boldsymbol{A},\boldsymbol{d}\right)$
\item $\bar{\mathcal{D}}=\left(\bar{\mathcal{N}},\bar{\mathcal{A}}\right)$ = directed graph in an atomic instance
\item $\left(\bar{S},\bar{I},\bar{T}\right)$ = (sources, pools, sinks) in an atomic instance
\item $\left(J_2,\ldots,J_6\right)$ = partition of $I$
\item $\boldsymbol{B}=\left(B_2,\ldots,B_6\right)$ = total capacities of pools in $J_2,\ldots,J_6$
\item $\zeta\left(\mathcal{D},\boldsymbol{b}\right)$ = minimum cost of flow in network $\mathcal{D}$ with node capacities $\boldsymbol{b}$ (dependency on other parameters suppressed)
\item $\left(e_2,\ldots,e_6\right)$ = unit vector representing the choice of one of the pool sets $J_2,\ldots,J_6$
\item $E\left(\boldsymbol{B}\right)$ = set of feasible values of $\left(e_2,\ldots,e_6\right)$ when $J_2,\ldots,J_6$ have capacities $B_2,\ldots,B_6$
\item $z_n\left(\boldsymbol{B}\right)$ = minimum cost when $J_2,\ldots,J_6$ have capacities $B_2,\ldots,B_6$ and $n-1$ pools are assigned to a set in $\left(J_2,\ldots,J_6\right)$
\item $\nu\left(\mathcal{D},\boldsymbol{b},\boldsymbol{c},\boldsymbol{q},\boldsymbol{u}\right)$ = size of a problem instance
\item $p$ = polynomial function of the instance size 
\end{description}
\end{document}

